% =============================================================================
% =============================================================================
% Which Derivatives Should You Label When the Simulator Is Wrong?
% NeurIPS 2026 workshop submission (ANONYMOUS). Target: Sim2Science, Paris.
% Deadline 2 Sep 2026 AoE. Reviews due 23 Sep.
%
% =============================================================================
% WHAT THIS DRAFT IS  (v5 -- label-budget reframe)
% =============================================================================
% Same experiments, same numbers, same frozen pre-registration. The headroom
% framing is GONE from the headline. Headroom was the gap to an oracle: real,
% but elementary and not novel, and three previous drafts were built on it.
%
% THE CONTRIBUTION IS NOW THE LABEL BUDGET. Three findings, none of which is
% in the prior work, all of which are already in your tables:
%   (A) Derivative supervision transfers off-distribution: 82-93% RMSE
%       reduction on two IN-CLASS coefficient regimes excised from training.
%       NOT a model-class departure -- Bates/Merton are hedging-PnL-only by
%       contract (splits.ood_models_hedging_only), there is no field metric
%       there, and at zero cost the value-only arm captures >=96% of the
%       downstream gap. Do not restore the old "model-class" wording.
%   (B) The CURVATURE rung carries the gain, and it TRANSFERS. Rung 2 adds only
%       d2u/dS2 labels and improves du/dv from 2.66/1.73 to 1.59/1.28 -- a
%       derivative it never supervises. This was in Table 1 and unremarked in
%       every previous draft.
%   (C) MORE LABELS IS NOT MONOTONICALLY BETTER. Rung 3 beats rung 2 on du/dv
%       but LOSES on d2u/dS2 (0.82 vs 0.64) and du/dS (0.0038 vs 0.0034) at the
%       extrapolated regime; registered dose-response monotonicity is not
%       demonstrated (rho_S = 0.029); and rung 2 vs rung 1 downstream is a
%       registered REVERSAL, harmful in sign though below the negligibility
%       floor. A ladder is the wrong mental model; it is a budget with a
%       trade-off.
% Section 3 is (A),(B),(C) then the controls then the downstream coda. DO NOT
% REORDER. The downstream null is a 200-word coda, which is its correct size.
%
% =============================================================================
% ONE INCONSISTENCY IN YOUR OWN SOURCES -- RECONCILE BEFORE SUBMITTING
% =============================================================================
% paper_s2s.tex says supervision cuts second-derivative and second-state RMSE
% "by 78 to 96%"; the headroom draft says "82 to 93%". The reduction rows of
% the RMSE table give 89/86 and 90/89 for d2u/dS2 and 93/93 and 87/82 for
% du/dv, i.e. a range of 82 to 93. This draft uses 82 to 93 because that is
% what the table supports. If 78-96 is right, the table is wrong; fix one.
%
% =============================================================================
% CONCESSIONS THAT MUST STAY (each closes an attack that is fatal alone)
% =============================================================================
%   - Two of the four reference solvers (Carr-Madan FFT, Fang-Oosterlee COS)
%     share a characteristic-function core and fail together where the branch
%     cut bites, which is the extrapolated-correlation anchor. Now cited to
%     albrecher2007little, which names that failure exactly.
%   - The 1e-3 agreement certifies the SOLUTION. Section 2 points to
%     app:oracle for the per-derivative spread. THAT APPENDIX MUST REPORT IT.
%   - Masking removes disagreeing points, which concentrate in the regions the
%     appendix names.
%   - The residual arm uses a fixed weight in the style of the cited PINN, not
%     a tuned one; adaptive loss balancing was not run.
%   - Two of three controls land ABOVE baseline (0.982, 0.699).
%   - Labels here are NOISELESS; coarse-grained force matching's are not.
%   - The registered accuracy endpoint failed on a two-sided value-parity
%     guard, not on accuracy. Stated plainly.
%
% NO FINANCE VOCABULARY IN THE BODY. J is defined once as the mean of the worst
% 5% of outcomes. Do not reintroduce CVaR, hedging, rebalancing, Greeks or
% options. The one exception is the disclosure sentence in Section 2.
%
% =============================================================================
% PAGE BUDGET -- COMPILE AND CHECK
% =============================================================================
% ~2500 words, 1 table, 0 figures. The f0_problem figure was REMOVED to pay for
% the reviewer-response edits; its numbers now appear as a clause in Section 1.
% COMPILE AND CONFIRM. If you come in comfortably under 4.7 pp, put the figure
% back at 0.55\linewidth after the first paragraph of Section 1 and cut that
% clause. If you run OVER, cut in this order:
%   (a) Section 3.6's second paragraph (the grid/occupancy caution). (~5 lines)
%   (b) The loss-decomposition sentence in 3.5.                      (~4 lines)
%   (c) Table 1's two "reduction vs." rows; state them in prose.     (~3 lines)
% Do NOT cut: the transfer result in 3.3, the non-monotonicity paragraph in
% 3.4, the "How wrong is the world" paragraph in 3.1, the evidence-base
% paragraph before 3.2, the noiseless-labels limitation, or the parity-guard
% disclosure. Those are the reviewer-response edits and each one closes an
% objection that is fatal on its own.
%
% =============================================================================
% v4 CHANGELOG -- claim-narrowing pass (body net +22 lines vs v3)
% =============================================================================
% The ladder is NESTED. That single fact drove most of this pass: curvature
% enters at the same moment as rung 2, so "curvature is special" and "rung 2 is
% special" are confounded and cannot be separated without S+v or SS-alone.
%   - TITLE changed (revert note sits above \title).
%   - ABSTRACT splits survival (out-of-class anchors) from budget (in-class
%     held-out regimes), and states we correct no simulator.
%   - NEW intro para "What this is not": scope vs. the Sim2Science CFP. We
%     reduce surrogate approximation error, we do not mitigate sim-to-real
%     discrepancy. This is the venue-fit fix.
%   - NEW para in 3.3 "What the nesting cannot separate", naming S+v and SS.
%   - 3.2 retitled; "attributable to" -> "largest increment in this ordering".
%   - "reallocates capacity" -> "trades accuracy" (mechanism not demonstrated).
%   - Universal recommendations removed from 3.3, 3.4 and 5.
%   - "severity sweep" -> "sweep over the departure parameters" in 3.1; the
%     registered name is now flagged as overstated where it appears.
%   - 3.6 coda leads with the failure standing, not with the accounting.
%   - FIXED dangling \ref{sec:readout} -> sec:controls. VERIFY THIS IS RIGHT.
%   - Paid for with the author's own cut list: Table 1's two reduction rows
%     (folded into 3.1 prose, incl. the 91-97% first-derivative figures) and
%     the loss-decomposition sentence in 3.5. Cut (a), the 3.6 grid/occupancy
%     caution, was NOT taken -- it is the only off-policy warning and cutting
%     it orphans ross2011reduction.
%
% =============================================================================
% v5 APPENDIX PASS -- what changed, and the three things you must read
% =============================================================================
% All 11 \FILL slots are filled from the frozen artifacts. ./check_submission.sh
% passes: no live \FILL, 24 refs resolved, 31 citations resolved, environments
% balanced, checklist wired. It does NOT compile -- there is no LaTeX toolchain
% on this machine, so the PAGE COUNT IS STILL UNVERIFIED. Compile against the
% real neurips_2026.sty (it IS the real one, 443 lines, the "stub" note in v4
% was stale) before uploading.
%
% ---------------------------------------------------------------------------
% 1. THE HIGH-RISK CHECK CLEARED, BUT NOT FOR THE EXPECTED REASON
% ---------------------------------------------------------------------------
% The worry was that 3.4's non-monotonicity (0.641->0.823, 0.003399->0.003784 at
% the EXTRAPOLATED anchor) sits inside the reference spread because the
% Carr-Madan/COS branch cut bites at rho=-0.9. Measured: those gaps are 0.09 and
% 0.97 envelopes, so yes, inside -- but the branch cut has nothing to do with it.
% Recomputing the deterministic legs on the frozen grids, CF and FD agree to
% 1.6e-8 (first derivative) and 3.3e-8 (curvature) RMS there. The envelope is
% MONTE-CARLO NOISE: the stochastic leg sets it on 96% of retained points. The
% 2.65x envelope inversion DECOMPOSES (from the frozen mc_se, no resampling):
% 2.00 path budget (the 4x multiplier near_feller draws, exact 1/sqrt(n)) x 1.32
% regime (curvature is genuinely noisier per path at rho=-0.9), and 2.00*1.32
% reproduces 2.65 to 1%. So it is MOSTLY but NOT ONLY the budget -- an earlier
% draft of this note overclaimed. At matched budget delta and vega ratios are
% 0.84 and 0.70, i.e. the extrapolated anchor is BETTER resolved on those. 3.4 was ALREADY written correctly (it says
% "cannot resolve", not "degrades"); app:oracle now proves it and says why.
% Sharper: at that anchor rungs 1/2/3 have curvature RMSE 0.77/0.31/0.40 times
% the envelope -- ALL of them beat the reference's own certainty, so NO ordering
% among them is resolvable there. Only the baseline gap is.
%
% ---------------------------------------------------------------------------
% 2. ONE CONTROL IS NOT CONFIGURATION-MATCHED -- 3.5 WAS REWRITTEN
% ---------------------------------------------------------------------------
% The information-matched (extra-value-data) arm trained at lambda_pde = 1.0,
% 100x the shared 1e-2 and the WORST of the four registered candidates on the
% study's own selection score (0.372 vs 0.170). Verified on all five
% checkpoints; results/grid/README.md documents it; run_info_matching.py never
% reads lambdas_selected.yaml. Every other arm is at 0.01, or legitimately 0.
% CONSEQUENCE: "extra value data lands ABOVE baseline (0.699 vs 0.410)" is NOT
% supported -- at one budget unit and the selected weight that arm simply IS the
% baseline, so an unknown share of the gap is the weight. What survives is ONLY
% the sweep read against ITSELF, and only AT THAT WEIGHT: 5x data moves validation
% 0.647->0.596, plateau at 0.605, 2x width 0.709, none of it near supervision's
% 0.045. The general claim "value data cannot substitute for derivative labels"
% is NOT established either -- the matched-weight (0.01) budget sweep was never
% run. That rerun is the cheapest missing experiment in the paper. 3.5 and
% app:exhibits now say exactly this. The old "two of three controls land above
% baseline" line is now ONE (0.982).
%
% ---------------------------------------------------------------------------
% 3. A BETTER FRAMING IS AVAILABLE AND IS PARTLY TAKEN
% ---------------------------------------------------------------------------
% 3.1 used to say we "did not make" a severity-versus-advantage curve. We did
% make an ORDINAL one and it is clean: at zero cost the advantage over the
% baseline falls monotonically in EACH departure coordinate (all six chains),
% +27.7% in class -> +20.5/+17.9/+15.6 -> +10.6/+7.0/+4.9 -> +0.8 -> -3.0/-4.2,
% every pooled CI excluding zero except the sign-crossing cell. That is a
% measured decay to a sign change, not "no favourable regime". The registered
% no_decisive_regime verdict reads the 1% tier ONLY, where everything is harm.
% 3.1 now states the ordering; app:config Table 6 is the table. It stays INSIDE
% "How wrong is the world", keeps the "the name overstates it" concession, and
% ends by saying this is an ordering and not the curve a distance would buy.
% Title, abstract and section order are UNTOUCHED.
% DO NOT promote it to the headline for this submission. The reason is not
% modesty: the ordering is ordinal and descriptive, the registered readout is
% the 1% tier, and there the sweep is uniformly harmful. "Downstream value is
% conditional" is the defensible headline and it is the one in the title. The
% severity decay is a finding to build the NEXT paper on -- with a calibrated
% departure axis, and a cost tier that is not the one that inverts every arm.
%
% ---------------------------------------------------------------------------
% OTHER CORRECTIONS MADE TO THE BODY (all were checkable errors)
% ---------------------------------------------------------------------------
%  - 3.1 "Against the residual-only baseline ... 82 to 93% ... 91 to 95%": the
%    82% comes from the VALUE-ONLY baseline, not the residual-only one. Now
%    "against the two unsupervised baselines ... 82 to 93% and 91 to 97%".
%  - 3.4 "per-seed range [+0.0311,+0.0856]": that is the ACROSS-SEED 95% t
%    interval on the mean. The per-seed values span [-0.0072,+0.1181] and one
%    of ten seeds favours rung 2. Now stated correctly.
%  - Section 2 masking: the mask does NOT concentrate at the extrapolated
%    anchor. That anchor is the LESS masked of the two (9.2% vs 17.1%). It
%    concentrates at tau=0.04 (94%/79%) and JUST out of the money (47.5% in
%    S/K in [0.75,0.9)), not deep OTM.
%  - 3.6 rank exchange "in all five seeds": true 5/5 for the FIRST derivative
%    at both regimes; curvature is 3/5 and 4/5 on the grid side. Narrowed.
%  - Amendment dates in app:thresholds verified against git (2026-07-31,
%    08-04, 08-05, 08-05); all precede the 08-12 label freeze and the 08-18
%    results, so "Pre-eval: yes" holds.
%
% ---------------------------------------------------------------------------
% STILL OPEN
% ---------------------------------------------------------------------------
%  - COMPILE and check the page count. Appendices added ~4 pages of tables;
%    if the workshop counts appendices toward the limit, cut app:config's
%    Table 4 (per-arm compute) and app:exhibits' Table 9 (region) first.
%  - Nominate the reciprocal reviewer: a co-author with a publication record,
%    not the student first author. Two reviews by 23 Sep.
%
% =============================================================================
% PAGE BUDGET -- COMPILE AND CHECK
% =============================================================================
% Body ~2500 words, 1 table, 0 figures. If the BODY runs over 4.7 pp, cut in
% this order: (a) 3.6's second paragraph (grid/occupancy caution); (b) the
% loss-decomposition sentence in 3.5; (c) Table 1's reduction rows.
% Do NOT cut: the transfer result in 3.3, the non-monotonicity paragraph in
% 3.4, "How wrong is the world" in 3.1, the evidence-base paragraph before 3.2,
% the noiseless-labels limitation, the parity-guard disclosure, or the
% lambda_pde disclosure now in 3.5. Each closes an objection fatal on its own.
%
\documentclass{article}
 
\usepackage[dblblindworkshop,final]{neurips_2026}
 
\usepackage[utf8]{inputenc}
\usepackage[T1]{fontenc}
\usepackage{hyperref}
\usepackage{url}
\usepackage{booktabs}
\usepackage{amsfonts}
\usepackage{nicefrac}
\usepackage{microtype}
\usepackage{xcolor}
\usepackage{amsmath, amssymb}
\usepackage{graphicx}
\usepackage{listings}
\graphicspath{{figures/}}
 
\lstset{extendedchars=true, inputencoding=utf8,
        literate={---}{{---}}1 {--}{{--}}1}

% -----------------------------------------------------------------------------
% \FILL renders in red and is deliberately ugly. Every appendix slot below that
% needs your real numbers uses it. Nothing marked \FILL may survive to the
% submitted PDF. Run ./check_submission.sh before uploading; it greps for these
% and fails the build if any remain.
% -----------------------------------------------------------------------------
\newcommand{\FILL}[1]{\textcolor{red}{\textbf{[FILL: #1]}}}
 
% TITLE CHANGED. The old title, "Which Derivatives Should You Label When the
% Simulator Is Wrong?", asserts two things the experiments do not establish:
% (i) that a general answer exists (the ladder is NESTED -- S+v and SS+v were
% never trained), and (ii) that the budget result was obtained under
% model-class departure (it was not; it comes from in-class held-out regimes).
% To revert, delete the two lines below and restore the one-line original.
\title{When the Simulator Is Wrong, Do Better Derivatives Improve Decisions?}

\workshoptitle{Sim2Science: ML with Imperfect Scientific Models}

% Camera-ready author block.
\author{%
  Melvin Loh \\
  Nanyang Technological University \\
  Turing AI Scholars Programme \\
  College of Computing and Data Science \\
  \texttt{mloh018@e.ntu.edu.sg} \\
  \And
  Prathamesh Dinesh Joshi \\
  Vizuara AI Labs \\
  \texttt{prathamesh@vizuara.com} \\
}

\hypersetup{hidelinks, pdfauthor={}, pdfsubject={}, pdfkeywords={}}

\begin{document}
\maketitle

\begin{abstract}
We use option pricing and hedging to study whether better surrogate derivatives
improve decisions when the assumed simulator is wrong. A neural Heston surrogate
predicts option values. A fixed hedging rule uses their sensitivity to asset price
to choose an asset holding intended to offset changes in the option's value.
At two Heston coefficient regimes excluded from training and validation,
adding curvature labels to price-sensitivity supervision reduces root-mean-square
error (RMSE) in the unsupervised variance sensitivity by $40$ and $26\%$.
Supervising all three targets reduces curvature and variance-sensitivity RMSE by
$82$--$93\%$ relative to two baselines without derivative labels. The incremental
gains depend on the target and are assessed against measured solver disagreement.
When simulated dynamics depart from the Heston model class, a value-only model
captures at least $96\%$ of each of the fully supervised model's eight improvements
over the residual-only baseline at zero transaction cost. At the registered $1\%$
cost, the fully supervised model is worse than that baseline in all ten departure
settings. The pre-registered confirmatory hedging endpoint, evaluated under a
coefficient shift, also fails. These findings support cross-target transfer within
the tested nested design, without establishing a consistent hedging advantage
specific to derivative labels or a general ordering of label choices.
\end{abstract}

% =============================================================================
\section{Introduction}
\label{sec:intro}

Scientific surrogates approximate models that can themselves be wrong. We use
option pricing and hedging to ask whether reducing derivative error improves
decisions under model misspecification. The network predicts a Heston option value;
a fixed hedging rule uses its asset-price sensitivity. Numerical references measure
value and derivative errors; simulations outside Heston test whether these errors
matter for hedging.
A close value fit does not ensure accurate derivatives: at a held-out coefficient
regime, two models without derivative labels have value errors of $1.6\%$ and
$3.3\%$, but curvature errors of $41\%$ and $27\%$
(Appendix~\ref{app:exhibits}). A physics residual also leaves substantial derivative
error, consistent with documented limitations
\citep{krishnapriyan2021characterizing,wang2021understanding}.

Derivative supervision is established in Sobolev training
\citep{czarnecki2017sobolev} and differential machine learning
\citep{huge2020differential}. Molecular models fit forces
\citep{ercolessi1994force,wang2019machine} or differentiate learned energies
\citep{behler2007generalized}. Sobolev-PINNs augment residual training
\citep{son2021sobolev}. Our contribution is a controlled evaluation of which
targets improve and whether this improves hedging. Model misspecification also
affects simulation-based inference \citep{cranmer2020frontier,cannon2022investigating}.
Sequential prediction also highlights the importance of evaluation distributions
\citep{ross2011reduction}.

\textbf{Does fidelity translate under misspecification?} We distinguish held-out
coefficients within the assumed model from a generating process outside its class,
which no parameter setting recovers. In the financial setting, this means asking
whether more accurate Heston derivatives improve hedging when the simulated price
process departs from Heston dynamics.

\textbf{Where should the labels go?} Derivative labels require additional reference
computations. We add targets sequentially to measure their conditional contribution.
\citet{noguer2026derivative} identifies the
importance of second derivatives in discrete-time hedging but reports no hedging
backtest. \citet{malandain2025deepsvm} reports inaccurate derivatives
despite a small residual and proposes derivative regularisation. The closest trained
configuration \citep{sakuma2026a} supervises first and second derivatives together,
but evaluates in distribution and reports similar pricing and Greek accuracy
among its three supervised arms.

We compare a nested sequence of label sets, pre-registered controls, and numerical
references with measured uncertainty. The registered tests and thresholds were fixed
before training. Numerical cross-checks let us measure errors in the value and the
evaluated derivatives, and assess whether differences exceed reference uncertainty.

\textbf{What we find.} Full supervision reduces held-out curvature and
variance-derivative RMSE by $82$--$93\%$ relative to the value-only and residual-only
baselines. Adding curvature labels also improves the
unlabelled variance derivative by $40$ and $26\%$, but adjacent-rung gains depend on
the target. Under model-class departure, value-only training largely reproduces
the zero-cost hedging gains; at $1\%$ cost the fully supervised model is worse than
the residual-only baseline in all ten settings. The registered confirmatory test
also fails (Sections~\ref{sec:survives} and~\ref{sec:coda}).

\textbf{What this is not.} We do not correct the simulator's dynamics. We
reduce the surrogate's error against an assumed pricing model and ask whether this
improves a fixed hedging rule when the model is structurally wrong. The experiments
test this distinction in one PDE setting; they do not establish transfer to other
simulators or real-market performance.

% =============================================================================
\section{Setup}
\label{sec:setup}

\textbf{Model and notation.} The Heston model \citep{heston1993closed} describes an
asset price $S$ and its instantaneous variance $v$, both evolving randomly. The
function $u(S,v,\tau;K)$ is the model value of a European call option: a contract
giving its holder the right to buy the asset at a fixed price $K$ (the strike)
at expiry. With time to expiry $\tau$, its payoff at expiry is $\max(S-K,0)$. For fixed
coefficients, the pricing PDE is
\begin{equation}
\label{eq:heston}
\begin{aligned}
\partial_\tau u ={}& \tfrac12 vS^2\partial_{SS}u
  + \rho\xi vS\partial_{Sv}u + \tfrac12\xi^2v\partial_{vv}u \\
  &+ (r-q)S\partial_Su + \kappa(\theta-v)\partial_vu - ru,
\end{aligned}
\end{equation}
with $u(S,v,0;K)=\max(S-K,0)$. Here $\kappa$ is the variance mean-reversion
rate, $\theta$ its long-run level, $\xi$ the variance diffusion scale, and $\rho$ the
correlation between the price and variance shocks; $r$ and $q$ are the interest rate
and dividend yield. The variance input is denoted $v_0$ in the data and configuration.
The network approximates $u$ across coefficients, as in parametric Heston PINNs
\citep{hainaut2024option}; the squared PDE residual measures its failure to satisfy
Eq.~\eqref{eq:heston} and supplies the physics-informed loss \citep{raissi2019physics}.

The three targets are \emph{delta}, the price sensitivity $\partial_Su$;
\emph{gamma}, the curvature $\partial_{SS}u=\partial^2u/\partial S^2$; and the
variance sensitivity $\partial_vu$. The last is a derivative with respect to
variance, rather than volatility. Delta describes how the option value changes
with asset price; gamma describes how that sensitivity itself changes. All are input derivatives obtained by automatic
differentiation, distinct from loss gradients with respect to network weights.
The registered hedging rule uses delta to set its asset holding; gamma and variance
sensitivity are supervision targets, not additional inputs to that rule.

\textbf{Reference cross-checks.} We cross-check option values using
up to four solvers: Heston Fourier inversion
\citep{heston1993closed,albrecher2007little}, cosine expansion
\citep{fang2008novel}, an alternating-direction finite-difference scheme
\citep{craig1988alternating}, and Monte Carlo \citep{lord2010comparison}.
The first two are separately implemented but share a characteristic-function
formula, so they cannot exclude a common formula error. The predeclared mask combines
agreement checks across values and derivatives, with a relative tolerance of
$10^{-3}$, a scale floor near zero, and an allowance for Monte-Carlo noise
(Appendix~\ref{app:oracle}). It removes $17.1\%$ of the interpolated grid
and $9.2\%$ of the extrapolated grid, concentrated at short expiries and just below
the strike. For each metric, all models are evaluated on the same retained points.

\textbf{What one envelope is.} Value agreement does not certify derivative accuracy. For each derivative and
evaluation regime, an \emph{uncertainty envelope} combines the largest pairwise
solver disagreement with Monte-Carlo standard error (Table~\ref{tab:spread}). A
difference is called \emph{resolved} when it exceeds that envelope; this is a
resolution heuristic, not a statistical test or error bound. Unresolved
differences need not be zero. Labels and reference share a
solver family: holding out parameters prevents memorisation of those regimes but
cannot exclude shared bias. Appendix~\ref{app:oracle} quantifies these limits.

\textbf{Departure.} An \emph{anchor} is a fixed coefficient setting
excluded, together with a neighbourhood, from both training and validation.
Derivative accuracy is evaluated at two Heston anchors: one inside the sampled
coefficient ranges and one with correlation outside them. These test interpolation
into a held-out region and extrapolation within the model class. For departures
outside that class, we evaluate hedging under Bates dynamics, which add price jumps
\citep{bates1996jumps}, and Merton dynamics, which combine jumps with constant
volatility \citep{merton1976option}. In these tests every model receives the same
assumed Heston coefficients; no derivative-accuracy claim is made for the departed
dynamics (Appendix~\ref{app:config}).

\textbf{Arms.} An \emph{arm} is one training configuration.
Both baselines use value labels: the value-only model has no PDE residual, whereas
the residual-only baseline adds that residual but no derivative labels. A nested
\emph{ladder} adds labels to the residual baseline: rung 1 adds delta, rung 2 also
adds gamma, and rung 3 also adds variance sensitivity. Each label loss is a
scale-normalised squared error \citep{huge2020differential}. These arms share
architecture, optimiser, data and budget. The residual, gamma and variance-label
weights were selected on validation; the delta-label weight was fixed. Adaptive
loss balancing was not tested. The extra-value-data control
uses a different residual weight, disclosed in Section~\ref{sec:controls}.

The hedging rule holds $\widehat{\partial_Su}$ units of the
asset to offset changes in the option's value. It updates that holding over $43$
steps on $10{,}000$ shared simulated trajectories
\citep{andersen2008simple}. Hedging loss $J$ is expected shortfall, the mean of the
worst $5\%$ of losses, so it measures adverse outcomes rather than average performance
\citep{rockafellar2000optimization}; lower is better.
Tests include zero and nonzero transaction costs. The pre-registered thresholds
require a $10\%$ confirmatory improvement, a $15\%$ derivative-accuracy improvement,
and specify a $2\%$ negligibility floor (Table~\ref{tab:verdicts}).

% =============================================================================
\section{Results}
\label{sec:results}

We report the registered intervals as specified before training. Later comparisons
are exploratory, without adjustment for multiple comparisons. Path-bootstrap
intervals measure variation across simulated paths for the fitted models, not
variation across new training runs.

\begin{table}[!htbp]
\caption{Held-out derivative RMSE on the full parameter grid, interpolated /
extrapolated coefficient regime, five seeds; second-derivative entries $\times
10^{-3}$. Lower is better. Bold marks a lowest mean separated from the runner-up by
more than the reference uncertainty; an unbolded regime--derivative cell is one whose
top two arms this benchmark cannot separate. Resolution is the registered
reference-noise heuristic, not a test. Lower block: those contrasts divided by the
reference envelope for that derivative; italics mark a gap under one envelope, which
this benchmark cannot resolve (Appendix~\ref{app:oracle}).}
\label{tab:rmse}
\centering\footnotesize\setlength{\tabcolsep}{6pt}
\begin{tabular}{@{}lrrr@{}}
\toprule
Arm & $\partial u/\partial S$ & $\partial^{2}u/\partial S^{2}$ & $\partial u/\partial v$ \\
\midrule
residual-only baseline & $0.0537/0.0424$ & $7.82/6.04$ & $9.93/7.44$ \\
value-only feedforward & $0.0881/0.0815$ & $8.29/7.17$ & $5.74/2.89$ \\
rung 1, $+\partial u/\partial S$ & $0.0106/0.0079$ & $2.56/1.58$ & $2.66/1.73$ \\
rung 2, $+\partial^{2}u/\partial S^{2}$ & $0.0033/0.0034$ & $1.02/0.64$ & $1.59/1.28$ \\
rung 3, $+\partial u/\partial v$ & $\mathbf{0.0027}/0.0038$ & $0.83/0.82$ & $\mathbf{0.73}/\mathbf{0.52}$ \\
\midrule
\multicolumn{4}{@{}l}{\emph{same contrasts, in reference envelopes}} \\
baseline $\to$ rung 3 & $184.7/97.3$ & $9.1/2.6$ & $85.0/54.5$ \\
rung 1 $\to$ rung 2 & $26.2/11.4$ & $2.0/\mathit{0.5}$ & $9.9/3.5$ \\
rung 2 $\to$ rung 3 & $2.2/\mathit{-0.97}$ & $\mathit{0.2}/\mathit{-0.1}$ & $8.0/6.0$ \\
\bottomrule
\end{tabular}
\end{table}

\subsection{Derivative accuracy and hedging under model departure}
\label{sec:survives}

Table~\ref{tab:rmse} evaluates derivative accuracy at held-out coefficients within
the Heston class. Relative to the two baselines without derivative labels, full
supervision reduces gamma and variance-sensitivity RMSE by $82$--$93\%$ and delta
RMSE by $91$--$97\%$. These baseline-to-rung-3 improvements exceed the reference
uncertainty at both regimes (Appendix~\ref{app:oracle}). The derivative-accuracy
findings below concern these two regimes; departures outside the model class are
evaluated through hedging loss.

At zero transaction cost, rung 3 has lower mean loss than the residual-only baseline
in seven of nine Bates jump settings and in the Merton setting. The value-only model
captures at least $96\%$ of each of these eight gains and outperforms rung 3 in four.
At the registered $1\%$ cost, rung 3 has higher loss than the residual-only baseline
in all ten settings. Thus the observed gains over that baseline do not establish a
consistent hedging advantage specific to derivative labels.

\textbf{Extent of misspecification.} The Bates sweep varies jump arrival rate
$\lambda_j\leq0.5$ and jump-size standard deviation $\sigma_j\leq0.15$
(Appendix~\ref{app:config}). These parameters order the tested departures, but we
did not measure distance from the Heston model class. At zero cost the relative gain
over the residual-only baseline decreases along each coordinate, from $+27.7\%$
at the in-class setting to $-4.2\%$ at the largest tested jumps, changing sign along
the $\lambda_j=0.5$ row. All six one-dimensional chains decrease in every one of
five seeds, for both absolute and relative gains ($30/30$ each). This monotonicity
analysis is exploratory and differs from the preregistered expectation of a
non-monotone pattern. The registered departure-sweep verdict uses the $1\%$ cost
tier, where no favourable regime is found.

The registered derivative-accuracy endpoint also fails, despite gamma and
variance-sensitivity reductions of $86.4$--$93.0\%$ exceeding its $15\%$ threshold.
It additionally requires value accuracy to remain within a two-sided parity band;
improving value error by $82$--$91\%$ violates that band. We retain the registered
failure and disclose this limitation of the criterion.

\subsection{Adding curvature labels improves curvature accuracy}
\label{sec:curvature}

Adding gamma labels to delta supervision reduces mean gamma RMSE by $60\%$ and
$59\%$, while value error changes by $4.4\%$ at both regimes. The relative
curvature improvement is therefore much larger than the change in value fit in
this comparison. The gamma improvement is $2.0$ reference envelopes at the
interpolated regime and $0.5$ at the extrapolated regime: it is resolved only at
the former. The larger baseline-to-rung-2 improvement is resolved at both. This
conditional comparison neither establishes a general preference for gamma labels
nor changes the registered parity failure.

\subsection{The gain is not confined to the supervised derivative}
\label{sec:transfer}

Rung 2 adds only gamma labels to rung 1, yet RMSE in the unsupervised variance
sensitivity $\partial_vu$ falls from $2.66/1.73$ to $1.59/1.28$, a $40\%$ and
$26\%$ reduction. Delta error also falls two- to three-fold. The variance-sensitivity
gaps are $9.9$ and $3.5$ reference envelopes, so cross-target transfer is resolved
at both regimes. One possible explanation is coupling through the PDE residual,
which contains both $\partial_{SS}u$ and $\partial_vu$. That mechanism remains
untested; varying the residual and label terms independently would help evaluate it.

\textbf{What the nested design establishes.} This comparison measures the effect
of adding gamma labels when delta labels and the shared residual are already
present. It does not separate their interactions or establish that gamma is the
best next target. Missing configurations include delta plus variance labels and
gamma labels alone. These are targeted follow-ups; a full factorial design would
require all combinations of the three label types. The transfer result therefore
applies to this sequence of added labels.

\subsection{The third target yields a resolved gain primarily on itself}
\label{sec:budget}

Adding variance-sensitivity labels in rung 3 reduces that target's RMSE from
$1.59/1.28$ to $0.73/0.52$, differences of $8.0$ and $6.0$ reference envelopes.
Delta also improves at the interpolated regime by $2.2$ envelopes, while the gamma
improvement is unresolved. At extrapolation, mean delta RMSE increases from
$0.003399$ to $0.003784$, and gamma RMSE in units of $10^{-3}$ from $0.641$ to
$0.823$, but the
differences are only $0.97$ and $0.09$ envelopes. We therefore cannot conclude that
either derivative degrades. All three supervised arms have gamma errors below its
reference uncertainty there (Appendix~\ref{app:oracle}). Only the newly supervised
target improves beyond the reference uncertainty at both regimes.

The registered label-noise experiment also fails to demonstrate a monotone
dose-response. In the confirmatory hedging comparison, adding gamma labels to delta
supervision increases loss by a statistically nonzero amount below the $2\%$
negligibility floor. These are distinct downstream tests, reported with their
original verdicts in Table~\ref{tab:verdicts}; they do not identify the mechanism
behind the derivative-accuracy results.

\subsection{Controls for label information}
\label{sec:controls}

The pre-registered controls use a curvature penalty without labels, shuffled
labels, or additional value data. Against the baseline's held-out gamma relative
RMSE of $0.410$, the penalty gives $0.982$ and shuffled labels $0.408$; neither
reproduces the supervised improvement. The penalty weight was not tuned.
The price-only sweep uses $4{,}096$--$20{,}480$ rows, fewer than the main arms'
$23{,}623$, and a residual weight $100$ times larger. Its validation gamma relative
RMSE changes from $0.647$ to $0.596$; doubling width gives $0.709\pm0.061$.
Because these training conditions differ, the sweep cannot establish whether extra
value labels can substitute for derivative labels (Appendix~\ref{app:exhibits}).

\subsection{Hedging performance and the registered endpoint}
\label{sec:coda}

The confirmatory hedging test fails: the relative improvement is $0.02\%$ against a
$10\%$ requirement. The paired $95\%$ interval for the \emph{loss difference},
$[-0.0242,+0.0242]$, includes zero. Together with the model-class departures in
Section~\ref{sec:survives}, this limits the downstream claim: large improvements
in derivative accuracy do not establish a consistent hedging advantage specific to
derivative supervision in the evaluated settings.

A separate, exploratory zero-cost comparison considers the same coefficient shift.
The fully supervised model's hedging loss is $0.7\%$ above that of the same rule using exact
Heston delta. Of its remaining gap to an analytic minimum-variance hedge
\citep{follmer1991hedging,hull2017optimal}, $96.7\%$ is associated with the choice
of hedging rule. Models trained with gamma labels reconstruct that hedge from their
own derivatives to within the previously specified $2\%$ tolerance; the closest
without gamma labels falls $3.3\%$ short (Appendix~\ref{app:exhibits}). These
comparisons are exploratory and leave the registered failure unchanged. At $1\%$
cost, even the residual-only baseline beats exact Heston delta, which is therefore
not a lower bound on loss for this task.

\textbf{Accuracy rankings depend on the evaluation region.} Restricting delta
RMSE to the hedged contract's region reverses the two baselines: at the interpolated
anchor, the residual-only model is better on the full grid ($0.0537$ versus
$0.0881$), but worse regionally ($0.0937$ versus $0.0409$). The reversal holds in
all five seeds at both anchors (Table~\ref{tab:region}). This exploratory metric
is not trajectory-weighted and does not establish a general predictor of hedging
performance.

% =============================================================================
\section{Conclusion}
\label{sec:limits}

Derivative supervision improves accuracy at held-out Heston coefficients.
Adding curvature labels also improves the unsupervised variance derivative.
These gains do not establish a consistent hedging advantage specific to derivative
labels under the tested departures and costs. The confirmatory test fails,
reinforcing the need to evaluate decisions alongside fidelity to the assumed model.

The study covers one two-state PDE, a small network, simulated paths and selected
departures, with no real-market validation. Derivative-accuracy comparisons use five seeds;
ladder hedging comparisons use ten. Labels and reference share a solver family,
so solver agreement cannot exclude common bias. Near zero variance, boundary
treatment may affect variance sensitivity. No operator-learning baseline or
adaptive residual weighting was tested. The fixed hedge need not be optimal at
each transaction cost \citep{whalley1997asymptotic}.

Broader validation should test other PDEs, label-noise levels and combinations of
derivative labels. A factorial design would help distinguish each label's effect
from its interactions with the others. Whether supervising energy curvature also
improves force accuracy in molecular models remains untested here. The practical
lesson is to measure both derivative error and decision quality, stating where
errors are evaluated, how uncertain the reference is, and which costs the decision incurs.

\clearpage
\begin{ack}
This work received no specific funding. The authors declare no competing interests.
\end{ack}

\begin{thebibliography}{99}

\bibitem[Albrecher et~al.(2007)]{albrecher2007little}
Albrecher, H., Mayer, P.~A., Schoutens, W., and Tistaert, J. (2007).
\newblock The little Heston trap.
\newblock \emph{Wilmott}, 2007(1):83--92.
 
\bibitem[Andersen(2008)]{andersen2008simple}
Andersen, L. (2008).
\newblock Simple and efficient simulation of the Heston stochastic volatility model.
\newblock \emph{Journal of Computational Finance}, 11(3):1--42.
 
% VERIFIED 2026-09-01
\bibitem[Bates(1996)]{bates1996jumps}
Bates, D.~S. (1996).
\newblock Jumps and stochastic volatility: exchange rate processes implicit in
Deutsche Mark options.
\newblock \emph{The Review of Financial Studies}, 9(1):69--107.
 
% VERIFIED 2026-09-01 -- ADDED FOR SIM2SCIENCE
\bibitem[Behler and Parrinello(2007)]{behler2007generalized}
Behler, J. and Parrinello, M. (2007).
\newblock Generalized neural-network representation of high-dimensional
potential-energy surfaces.
\newblock \emph{Physical Review Letters}, 98(14):146401.
 
% VERIFIED 2026-09-01 -- ADDED FOR SIM2SCIENCE
\bibitem[Cannon et~al.(2022)]{cannon2022investigating}
Cannon, P., Ward, D., and Schmon, S.~M. (2022).
\newblock Investigating the impact of model misspecification in neural
simulation-based inference.
\newblock arXiv:2209.01845.

 
% VERIFIED 2026-09-01 -- ADDED FOR SIM2SCIENCE
\bibitem[Craig and Sneyd(1988)]{craig1988alternating}
Craig, I.~J.~D. and Sneyd, A.~D. (1988).
\newblock An alternating-direction implicit scheme for parabolic equations with
mixed derivatives.
\newblock \emph{Computers \& Mathematics with Applications}, 16(4):341--350.

\bibitem[Cranmer et~al.(2020)]{cranmer2020frontier}
Cranmer, K., Brehmer, J., and Louppe, G. (2020).
\newblock The frontier of simulation-based inference.
\newblock \emph{Proceedings of the National Academy of Sciences},
117(48):30055--30062.
 
\bibitem[Czarnecki et~al.(2017)]{czarnecki2017sobolev}
Czarnecki, W.~M., Osindero, S., Jaderberg, M., \'Swirszcz, G., and Pascanu, R. (2017).
\newblock Sobolev training for neural networks.
\newblock In \emph{Advances in Neural Information Processing Systems (NIPS)}, volume 30.

\bibitem[Ercolessi and Adams(1994)]{ercolessi1994force}
Ercolessi, F. and Adams, J.~B. (1994).
\newblock Interatomic potentials from first-principles calculations: The
force-matching method.
\newblock \emph{Europhysics Letters}, 26(8):583--588.

\bibitem[Fang and Oosterlee(2008)]{fang2008novel}
Fang, F. and Oosterlee, C.~W. (2008).
\newblock A novel pricing method for European options based on Fourier-cosine
series expansions.
\newblock \emph{SIAM Journal on Scientific Computing}, 31(2):826--848.
 
\bibitem[F\"ollmer and Schweizer(1991)]{follmer1991hedging}
F\"ollmer, H. and Schweizer, M. (1991).
\newblock Hedging of contingent claims under incomplete information.
\newblock In Davis, M.~H.~A. and Elliott, R.~J., editors, \emph{Applied Stochastic
Analysis}, volume~5 of \emph{Stochastics Monographs}, pages 389--414. Gordon and
Breach.
 
\bibitem[Hainaut and Casas(2024)]{hainaut2024option}
Hainaut, D. and Casas, A. (2024).
\newblock Option pricing in the Heston model with physics inspired neural networks.
\newblock \emph{Annals of Finance}, 20(3):353--376.
 
\bibitem[Heston(1993)]{heston1993closed}
Heston, S.~L. (1993).
\newblock A closed-form solution for options with stochastic volatility with
applications to bond and currency options.
\newblock \emph{The Review of Financial Studies}, 6(2):327--343.
 
\bibitem[Huge and Savine(2020)]{huge2020differential}
Huge, B. and Savine, A. (2020).
\newblock Differential machine learning.
\newblock arXiv:2005.02347.
 
% VERIFIED 2026-09-01
\bibitem[Hull and White(2017)]{hull2017optimal}
Hull, J. and White, A. (2017).
\newblock Optimal delta hedging for options.
\newblock \emph{Journal of Banking \& Finance}, 82:180--190.
 
\bibitem[Krishnapriyan et~al.(2021)]{krishnapriyan2021characterizing}
Krishnapriyan, A.~S., Gholami, A., Zhe, S., Kirby, R.~M., and Mahoney, M.~W. (2021).
\newblock Characterizing possible failure modes in physics-informed neural networks.
\newblock In \emph{Advances in Neural Information Processing Systems (NeurIPS)},
volume~34, pages 26548--26560.

 

\bibitem[Lord et~al.(2010)]{lord2010comparison}
Lord, R., Koekkoek, R., and van Dijk, D. (2010).
\newblock A comparison of biased simulation schemes for stochastic volatility
models.
\newblock \emph{Quantitative Finance}, 10(2):177--194.

\bibitem[Malandain et~al.(2025)]{malandain2025deepsvm}
Malandain, K.~A., Kalici, S., and Chakhoyan, H. (2025).
\newblock DeepSVM: Learning stochastic volatility models with physics-informed
deep operator networks.
\newblock arXiv:2512.07162.
 
% VERIFIED 2026-09-01
\bibitem[Merton(1976)]{merton1976option}
Merton, R.~C. (1976).
\newblock Option pricing when underlying stock returns are discontinuous.
\newblock \emph{Journal of Financial Economics}, 3(1--2):125--144.
 
 
\bibitem[Noguer i Alonso(2026)]{noguer2026derivative}
Noguer i Alonso, M. (2026).
\newblock Derivative-informed operator learning for finance: On-the-fly Greeks,
surfaces, hedging, and control.
\newblock arXiv:2606.05900.
 
\bibitem[Raissi et~al.(2019)]{raissi2019physics}
Raissi, M., Perdikaris, P., and Karniadakis, G.~E. (2019).
\newblock Physics-informed neural networks: A deep learning framework for solving
forward and inverse problems involving nonlinear partial differential equations.
\newblock \emph{Journal of Computational Physics}, 378:686--707.
 
\bibitem[Rockafellar and Uryasev(2000)]{rockafellar2000optimization}
Rockafellar, R.~T. and Uryasev, S. (2000).
\newblock Optimization of conditional value-at-risk.
\newblock \emph{Journal of Risk}, 2(3):21--41.
 
% VERIFIED 2026-09-01
\bibitem[Ross et~al.(2011)]{ross2011reduction}
Ross, S., Gordon, G., and Bagnell, D. (2011).
\newblock A reduction of imitation learning and structured prediction to no-regret
online learning.
\newblock In \emph{Proceedings of the 14th International Conference on Artificial
Intelligence and Statistics (AISTATS)}, pages 627--635.
 
\bibitem[Sakuma(2026)]{sakuma2026a}
Sakuma, T. (2026).
\newblock Differential machine learning for 0DTE options with stochastic
volatility and jumps.
\newblock arXiv:2603.07600.
 
\bibitem[Son et~al.(2023)]{son2021sobolev}
Son, H., Jang, J.~W., Han, W.~J., and Hwang, H.~J. (2023).
\newblock Sobolev training for physics-informed neural networks.
\newblock \emph{Communications in Mathematical Sciences}, 21(6):1679--1705.
 
% VERIFIED 2026-09-01 -- ADDED FOR SIM2SCIENCE
\bibitem[Wang et~al.(2019)]{wang2019machine}
Wang, J., Olsson, S., Wehmeyer, C., P\'erez, A., Charron, N.~E., de~Fabritiis, G.,
No\'e, F., and Clementi, C. (2019).
\newblock Machine learning of coarse-grained molecular dynamics force fields.
\newblock \emph{ACS Central Science}, 5(5):755--767.
 
\bibitem[Wang et~al.(2021)]{wang2021understanding}
Wang, S., Teng, Y., and Perdikaris, P. (2021).
\newblock Understanding and mitigating gradient flow pathologies in
physics-informed neural networks.
\newblock \emph{SIAM Journal on Scientific Computing}, 43(5):A3055--A3081.
 
\bibitem[Whalley and Wilmott(1997)]{whalley1997asymptotic}
Whalley, A.~E. and Wilmott, P. (1997).
\newblock An asymptotic analysis of an optimal hedging model for option pricing
with transaction costs.
\newblock \emph{Mathematical Finance}, 7(3):307--324.
 
\end{thebibliography}
\section*{Declaration of AI-assisted technologies}
The authors used generative AI and AI-assisted tools for manuscript drafting,
language refinement, and code drafting and debugging. The authors directed the
project, reviewed and verified the manuscript, the code, and all reported results,
and take full responsibility for the final manuscript.

% Sim2Science requires the NeurIPS checklist after references and before appendices.
\section*{NeurIPS Paper Checklist}

\begin{enumerate}

\item {\bf Claims}
    \item[] Question: Do the main claims made in the abstract and introduction accurately reflect the paper's contributions and scope?
    \item[] Answer: \answerYes{}
    \item[] Justification: The abstract and Section~\ref{sec:intro} separate the two evidence bases the paper rests on --- held-out coefficient regimes inside the model class for the field results, and the downstream cells for model-class departure --- state that the registered confirmatory endpoint failed, and scope every claim to this testbed and this nested label ordering. Section~\ref{sec:limits} states the principal limitations.
    \item[] Guidelines:
    \begin{itemize}
        \item The answer \answerNA{} means that the abstract and introduction do not include the claims made in the paper.
        \item The abstract and/or introduction should clearly state the claims made, including the contributions made in the paper and important assumptions and limitations. A \answerNo{} or \answerNA{} answer to this question will not be perceived well by the reviewers. 
        \item The claims made should match theoretical and experimental results, and reflect how much the results can be expected to generalize to other settings. 
        \item It is fine to include aspirational goals as motivation as long as it is clear that these goals are not attained by the paper. 
    \end{itemize}

\item {\bf Limitations}
    \item[] Question: Does the paper discuss the limitations of the work performed by the authors?
    \item[] Answer: \answerYes{}
    \item[] Justification: Section~\ref{sec:limits} discusses the single PDE and small network, simulated rather than real-market data, seed counts, shared solver bias, boundary treatment, missing baselines, and the fixed hedging rule. Section~\ref{sec:transfer} explains why the nested design cannot isolate all label effects and interactions. Section~\ref{sec:controls} discloses the unmatched control, and Appendix~\ref{app:oracle} reports reference uncertainty.
    \item[] Guidelines:
    \begin{itemize}
        \item The answer \answerNA{} means that the paper has no limitation while the answer \answerNo{} means that the paper has limitations, but those are not discussed in the paper. 
        \item The authors are encouraged to create a separate ``Limitations'' section in their paper.
        \item The paper should point out any strong assumptions and how robust the results are to violations of these assumptions (e.g., independence assumptions, noiseless settings, model well-specification, asymptotic approximations only holding locally). The authors should reflect on how these assumptions might be violated in practice and what the implications would be.
        \item The authors should reflect on the scope of the claims made, e.g., if the approach was only tested on a few datasets or with a few runs. In general, empirical results often depend on implicit assumptions, which should be articulated.
        \item The authors should reflect on the factors that influence the performance of the approach. For example, a facial recognition algorithm may perform poorly when image resolution is low or images are taken in low lighting. Or a speech-to-text system might not be used reliably to provide closed captions for online lectures because it fails to handle technical jargon.
        \item The authors should discuss the computational efficiency of the proposed algorithms and how they scale with dataset size.
        \item If applicable, the authors should discuss possible limitations of their approach to address problems of privacy and fairness.
        \item While the authors might fear that complete honesty about limitations might be used by reviewers as grounds for rejection, a worse outcome might be that reviewers discover limitations that aren't acknowledged in the paper. The authors should use their best judgment and recognize that individual actions in favor of transparency play an important role in developing norms that preserve the integrity of the community. Reviewers will be specifically instructed to not penalize honesty concerning limitations.
    \end{itemize}

\item {\bf Theory assumptions and proofs}
    \item[] Question: For each theoretical result, does the paper provide the full set of assumptions and a complete (and correct) proof?
    \item[] Answer: \answerNA{}
    \item[] Justification: The paper introduces no theorem or new formal proof. The governing equation and the variance-minimising control ratio used in the experiments are stated with their assumptions and cited in Section~\ref{sec:setup} and Appendix~\ref{app:exhibits}.
    \item[] Guidelines:
    \begin{itemize}
        \item The answer \answerNA{} means that the paper does not include theoretical results. 
        \item All the theorems, formulas, and proofs in the paper should be numbered and cross-referenced.
        \item All assumptions should be clearly stated or referenced in the statement of any theorems.
        \item The proofs can either appear in the main paper or the supplemental material, but if they appear in the supplemental material, the authors are encouraged to provide a short proof sketch to provide intuition. 
        \item Inversely, any informal proof provided in the core of the paper should be complemented by formal proofs provided in appendix or supplemental material.
        \item Theorems and Lemmas that the proof relies upon should be properly referenced. 
    \end{itemize}

    \item {\bf Experimental result reproducibility}
    \item[] Question: Does the paper fully disclose all the information needed to reproduce the main experimental results of the paper to the extent that it affects the main claims and/or conclusions of the paper (regardless of whether the code and data are provided or not)?
    \item[] Answer: \answerYes{}
    \item[] Justification: Section~\ref{sec:setup} and Appendices~\ref{app:thresholds}--\ref{app:exhibits} disclose the frozen acceptance block, coefficient sampling and anchor excision, the reference construction with its per-derivative uncertainty, architecture, objective, optimizer and schedule, seeds, common-random-number simulation, and the paired-bootstrap analysis needed to reproduce every reported comparison.
    \item[] Guidelines:
    \begin{itemize}
        \item The answer \answerNA{} means that the paper does not include experiments.
        \item If the paper includes experiments, a \answerNo{} answer to this question will not be perceived well by the reviewers: Making the paper reproducible is important, regardless of whether the code and data are provided or not.
        \item If the contribution is a dataset and\slash or model, the authors should describe the steps taken to make their results reproducible or verifiable. 
        \item Depending on the contribution, reproducibility can be accomplished in various ways. For example, if the contribution is a novel architecture, describing the architecture fully might suffice, or if the contribution is a specific model and empirical evaluation, it may be necessary to either make it possible for others to replicate the model with the same dataset, or provide access to the model. In general. releasing code and data is often one good way to accomplish this, but reproducibility can also be provided via detailed instructions for how to replicate the results, access to a hosted model (e.g., in the case of a large language model), releasing of a model checkpoint, or other means that are appropriate to the research performed.
        \item While NeurIPS does not require releasing code, the conference does require all submissions to provide some reasonable avenue for reproducibility, which may depend on the nature of the contribution. For example
        \begin{enumerate}
            \item If the contribution is primarily a new algorithm, the paper should make it clear how to reproduce that algorithm.
            \item If the contribution is primarily a new model architecture, the paper should describe the architecture clearly and fully.
            \item If the contribution is a new model (e.g., a large language model), then there should either be a way to access this model for reproducing the results or a way to reproduce the model (e.g., with an open-source dataset or instructions for how to construct the dataset).
            \item We recognize that reproducibility may be tricky in some cases, in which case authors are welcome to describe the particular way they provide for reproducibility. In the case of closed-source models, it may be that access to the model is limited in some way (e.g., to registered users), but it should be possible for other researchers to have some path to reproducing or verifying the results.
        \end{enumerate}
    \end{itemize}


\item {\bf Open access to data and code}
    \item[] Question: Does the paper provide open access to the data and code, with sufficient instructions to faithfully reproduce the main experimental results, as described in supplemental material?
    \item[] Answer: \answerNo{}
    \item[] Justification: No public code or data archive is linked in this version. The experimental configuration and evaluation procedures are described in the paper and appendices.
    \item[] Guidelines:
    \begin{itemize}
        \item The answer \answerNA{} means that paper does not include experiments requiring code.
        \item Please see the NeurIPS code and data submission guidelines (\url{https://neurips.cc/public/guides/CodeSubmissionPolicy}) for more details.
        \item While we encourage the release of code and data, we understand that this might not be possible, so \answerNo{} is an acceptable answer. Papers cannot be rejected simply for not including code, unless this is central to the contribution (e.g., for a new open-source benchmark).
        \item The instructions should contain the exact command and environment needed to run to reproduce the results. See the NeurIPS code and data submission guidelines (\url{https://neurips.cc/public/guides/CodeSubmissionPolicy}) for more details.
        \item The authors should provide instructions on data access and preparation, including how to access the raw data, preprocessed data, intermediate data, and generated data, etc.
        \item The authors should provide scripts to reproduce all experimental results for the new proposed method and baselines. If only a subset of experiments are reproducible, they should state which ones are omitted from the script and why.
        \item At submission time, to preserve anonymity, the authors should release anonymized versions (if applicable).
        \item Providing as much information as possible in supplemental material (appended to the paper) is recommended, but including URLs to data and code is permitted.
    \end{itemize}


\item {\bf Experimental setting/details}
    \item[] Question: Does the paper specify all the training and test details (e.g., data splits, hyperparameters, how they were chosen, type of optimizer) necessary to understand the results?
    \item[] Answer: \answerYes{}
    \item[] Justification: Section~\ref{sec:setup} gives the design and the validation-only selection rule; Tables~\ref{tab:trainconfig} and~\ref{tab:arms-s2s} give the architecture, optimizer and schedule, batch and collocation sizes, label artifact, seeds, and the per-arm label flags, together with the residual weight and price-point budget each arm actually carried, including the one control that departs from the shared configuration.
    \item[] Guidelines:
    \begin{itemize}
        \item The answer \answerNA{} means that the paper does not include experiments.
        \item The experimental setting should be presented in the core of the paper to a level of detail that is necessary to appreciate the results and make sense of them.
        \item The full details can be provided either with the code, in appendix, or as supplemental material.
    \end{itemize}

\item {\bf Experiment statistical significance}
    \item[] Question: Does the paper report error bars suitably and correctly defined or other appropriate information about the statistical significance of the experiments?
    \item[] Answer: \answerYes{}
    \item[] Justification: Section~\ref{sec:results} states that path-bootstrap intervals are conditional on the fitted models and that exploratory comparisons are unadjusted for multiplicity. Section~\ref{sec:coda} reports the confirmatory interval; Appendix~\ref{app:exhibits} gives the paired-bootstrap procedure and identifies point-estimate-only comparisons. Appendix~\ref{app:oracle} reports reference uncertainty and unresolved orderings.
    \item[] Guidelines:
    \begin{itemize}
        \item The answer \answerNA{} means that the paper does not include experiments.
        \item The authors should answer \answerYes{} if the results are accompanied by error bars, confidence intervals, or statistical significance tests, at least for the experiments that support the main claims of the paper.
        \item The factors of variability that the error bars are capturing should be clearly stated (for example, train/test split, initialization, random drawing of some parameter, or overall run with given experimental conditions).
        \item The method for calculating the error bars should be explained (closed form formula, call to a library function, bootstrap, etc.)
        \item The assumptions made should be given (e.g., Normally distributed errors).
        \item It should be clear whether the error bar is the standard deviation or the standard error of the mean.
        \item It is OK to report 1-sigma error bars, but one should state it. The authors should preferably report a 2-sigma error bar than state that they have a 96\% CI, if the hypothesis of Normality of errors is not verified.
        \item For asymmetric distributions, the authors should be careful not to show in tables or figures symmetric error bars that would yield results that are out of range (e.g., negative error rates).
        \item If error bars are reported in tables or plots, the authors should explain in the text how they were calculated and reference the corresponding figures or tables in the text.
    \end{itemize}

\item {\bf Experiments compute resources}
    \item[] Question: For each experiment, does the paper provide sufficient information on the computer resources (type of compute workers, memory, time of execution) needed to reproduce the experiments?
    \item[] Answer: \answerNo{}
    \item[] Justification: Table~\ref{tab:arms-s2s} reports per-arm wall clock, parameter count, derivative-evaluation count and peak process memory, and Table~\ref{tab:trainconfig} names the worker type as single-process CPU. The processor model and a project-wide compute total, including preliminary and discarded runs, were not recorded, so the paper does not meet this item's full accounting standard.
    \item[] Guidelines:
    \begin{itemize}
        \item The answer \answerNA{} means that the paper does not include experiments.
        \item The paper should indicate the type of compute workers CPU or GPU, internal cluster, or cloud provider, including relevant memory and storage.
        \item The paper should provide the amount of compute required for each of the individual experimental runs as well as estimate the total compute. 
        \item The paper should disclose whether the full research project required more compute than the experiments reported in the paper (e.g., preliminary or failed experiments that didn't make it into the paper). 
    \end{itemize}
    
\item {\bf Code of ethics}
    \item[] Question: Does the research conducted in the paper conform, in every respect, with the NeurIPS Code of Ethics \url{https://neurips.cc/public/EthicsGuidelines}?
    \item[] Answer: \answerYes{}
    \item[] Justification: This numerical study uses synthetic data, no human participants or private data, and reports failed registered endpoints and adverse diagnostics rather than suppressing them. AI-assisted drafting is disclosed before the appendices, and the authors retain responsibility for verification.
    \item[] Guidelines:
    \begin{itemize}
        \item The answer \answerNA{} means that the authors have not reviewed the NeurIPS Code of Ethics.
        \item If the authors answer \answerNo, they should explain the special circumstances that require a deviation from the Code of Ethics.
        \item The authors should make sure to preserve anonymity (e.g., if there is a special consideration due to laws or regulations in their jurisdiction).
    \end{itemize}


\item {\bf Broader impacts}
    \item[] Question: Does the paper discuss both potential positive societal impacts and negative societal impacts of the work performed?
    \item[] Answer: \answerNA{}
    \item[] Justification: This is a synthetic methodological evaluation of how derivative labels should be spent when training a surrogate, not a deployed system or a study involving people, so no direct societal impact is claimed. Its results should not be treated as validation of any decision system built on the testbed used here, a scope restriction stated in Sections~\ref{sec:intro} and~\ref{sec:limits}.
    \item[] Guidelines:
    \begin{itemize}
        \item The answer \answerNA{} means that there is no societal impact of the work performed.
        \item If the authors answer \answerNA{} or \answerNo, they should explain why their work has no societal impact or why the paper does not address societal impact.
        \item Examples of negative societal impacts include potential malicious or unintended uses (e.g., disinformation, generating fake profiles, surveillance), fairness considerations (e.g., deployment of technologies that could make decisions that unfairly impact specific groups), privacy considerations, and security considerations.
        \item The conference expects that many papers will be foundational research and not tied to particular applications, let alone deployments. However, if there is a direct path to any negative applications, the authors should point it out. For example, it is legitimate to point out that an improvement in the quality of generative models could be used to generate Deepfakes for disinformation. On the other hand, it is not needed to point out that a generic algorithm for optimizing neural networks could enable people to train models that generate Deepfakes faster.
        \item The authors should consider possible harms that could arise when the technology is being used as intended and functioning correctly, harms that could arise when the technology is being used as intended but gives incorrect results, and harms following from (intentional or unintentional) misuse of the technology.
        \item If there are negative societal impacts, the authors could also discuss possible mitigation strategies (e.g., gated release of models, providing defenses in addition to attacks, mechanisms for monitoring misuse, mechanisms to monitor how a system learns from feedback over time, improving the efficiency and accessibility of ML).
    \end{itemize}
    
\item {\bf Safeguards}
    \item[] Question: Does the paper describe safeguards that have been put in place for responsible release of data or models that have a high risk for misuse (e.g., pre-trained language models, image generators, or scraped datasets)?
    \item[] Answer: \answerNA{}
    \item[] Justification: The work introduces no pretrained model, scraped dataset, personal data, or other artifact with a high risk of misuse; it consists of small numerical surrogates and synthetically generated trajectories.
    \item[] Guidelines:
    \begin{itemize}
        \item The answer \answerNA{} means that the paper poses no such risks.
        \item Released models that have a high risk for misuse or dual-use should be released with necessary safeguards to allow for controlled use of the model, for example by requiring that users adhere to usage guidelines or restrictions to access the model or implementing safety filters. 
        \item Datasets that have been scraped from the Internet could pose safety risks. The authors should describe how they avoided releasing unsafe images.
        \item We recognize that providing effective safeguards is challenging, and many papers do not require this, but we encourage authors to take this into account and make a best faith effort.
    \end{itemize}

\item {\bf Licenses for existing assets}
    \item[] Question: Are the creators or original owners of assets (e.g., code, data, models), used in the paper, properly credited and are the license and terms of use explicitly mentioned and properly respected?
    \item[] Answer: \answerNA{}
    \item[] Justification: The experiments use no external dataset, pretrained model, or third-party research artifact. All labels and trajectories are generated from the equations and numerical methods specified and cited in the paper.
    \item[] Guidelines:
    \begin{itemize}
        \item The answer \answerNA{} means that the paper does not use existing assets.
        \item The authors should cite the original paper that produced the code package or dataset.
        \item The authors should state which version of the asset is used and, if possible, include a URL.
        \item The name of the license (e.g., CC-BY 4.0) should be included for each asset.
        \item For scraped data from a particular source (e.g., website), the copyright and terms of service of that source should be provided.
        \item If assets are released, the license, copyright information, and terms of use in the package should be provided. For popular datasets, \url{paperswithcode.com/datasets} has curated licenses for some datasets. Their licensing guide can help determine the license of a dataset.
        \item For existing datasets that are re-packaged, both the original license and the license of the derived asset (if it has changed) should be provided.
        \item If this information is not available online, the authors are encouraged to reach out to the asset's creators.
    \end{itemize}

\item {\bf New assets}
    \item[] Question: Are new assets introduced in the paper well documented and is the documentation provided alongside the assets?
    \item[] Answer: \answerNA{}
    \item[] Justification: No public asset release is linked in this version. The paper describes the generated labels and trajectories, training configurations, and evaluation procedures.
    \item[] Guidelines:
    \begin{itemize}
        \item The answer \answerNA{} means that the paper does not release new assets.
        \item Researchers should communicate the details of the dataset\slash code\slash model as part of their submissions via structured templates. This includes details about training, license, limitations, etc. 
        \item The paper should discuss whether and how consent was obtained from people whose asset is used.
        \item At submission time, remember to anonymize your assets (if applicable). You can either create an anonymized URL or include an anonymized zip file.
    \end{itemize}

\item {\bf Crowdsourcing and research with human subjects}
    \item[] Question: For crowdsourcing experiments and research with human subjects, does the paper include the full text of instructions given to participants and screenshots, if applicable, as well as details about compensation (if any)? 
    \item[] Answer: \answerNA{}
    \item[] Justification: The study involves neither crowdsourcing nor human subjects.
    \item[] Guidelines:
    \begin{itemize}
        \item The answer \answerNA{} means that the paper does not involve crowdsourcing nor research with human subjects.
        \item Including this information in the supplemental material is fine, but if the main contribution of the paper involves human subjects, then as much detail as possible should be included in the main paper. 
        \item According to the NeurIPS Code of Ethics, workers involved in data collection, curation, or other labor should be paid at least the minimum wage in the country of the data collector. 
    \end{itemize}

\item {\bf Institutional review board (IRB) approvals or equivalent for research with human subjects}
    \item[] Question: Does the paper describe potential risks incurred by study participants, whether such risks were disclosed to the subjects, and whether Institutional Review Board (IRB) approvals (or an equivalent approval/review based on the requirements of your country or institution) were obtained?
    \item[] Answer: \answerNA{}
    \item[] Justification: The study involves no human subjects, so institutional human-subjects review is not applicable.
    \item[] Guidelines:
    \begin{itemize}
        \item The answer \answerNA{} means that the paper does not involve crowdsourcing nor research with human subjects.
        \item Depending on the country in which research is conducted, IRB approval (or equivalent) may be required for any human subjects research. If you obtained IRB approval, you should clearly state this in the paper. 
        \item We recognize that the procedures for this may vary significantly between institutions and locations, and we expect authors to adhere to the NeurIPS Code of Ethics and the guidelines for their institution. 
        \item For initial submissions, do not include any information that would break anonymity (if applicable), such as the institution conducting the review.
    \end{itemize}

\item {\bf Declaration of LLM usage}
    \item[] Question: Does the paper describe the usage of LLMs if it is an important, original, or non-standard component of the core methods in this research? Note that if the LLM is used only for writing, editing, or formatting purposes and does \emph{not} impact the core methodology, scientific rigor, or originality of the research, declaration is not required.
    %this research? 
    \item[] Answer: \answerNA{}
    \item[] Justification: LLMs are not an important, original, or non-standard component of the core scientific method. The paper nevertheless voluntarily discloses their use for manuscript and code drafting and debugging in the declaration preceding the appendices; the authors reviewed and verified the work.
    \item[] Guidelines:
    \begin{itemize}
        \item The answer \answerNA{} means that the core method development in this research does not involve LLMs as any important, original, or non-standard components.
        \item Please refer to our LLM policy in the NeurIPS handbook for what should or should not be described.
    \end{itemize}

\end{enumerate}


\appendix
% =============================================================================
\section{Pre-registered acceptance block and amendment record}
\label{app:thresholds}
The decision-bearing portions below are transcribed from
\texttt{heston\_benchmark\_v6.yaml}, the contract consumed by the frozen label and
training artifacts. Long audit comments are omitted, but key names and values are
unchanged. The label artifact is frozen at tag
\texttt{v6-labels-20260812}.

\begin{lstlisting}[basicstyle=\ttfamily\tiny,breaklines=true,
  breakatwhitespace=true,keepspaces=true,columns=fullflexible]
acceptance_thresholds:
  confirmatory_cell_pass: "misspec delta CVaR95 improvement >=10% relative AND paired-bootstrap 95% CI excludes 0 (combined,1% TC,baseline,rung3 vs standard_pinn)"
  confirmatory_cell_rel_min: 0.10
  headline_scale_free: fraction_of_baseline_to_oracle_gap_closed
  order_attribution: "rung2 beats rung1 at the cell, CI excludes 0; failure = honest null"
  dose_response:
    criteria: "Spearman rho > 0 (shape confirmed via isotonic fit) AND one-sided seed-bootstrap tail probability P(rho<=0) < bootstrap_tail_prob_max"
    bootstrap_tail_prob_max: 0.05
    interpretation: "monotonicity not demonstrated is not the claim that the response is flat"
  ood_gamma_rmse_reduction_min: 0.15
  ood_vega_rmse_reduction_min: 0.15
  price_parity_within: 0.10
  in_model_hedging: NOT_PASS_FAIL
  sakuma_null_rel_tol: 0.02
  seeds: {default: 5, confirmatory_cell: 10}
  tail_claim_requires: paired_bootstrap_over_CRN_paths_with_seed_variance_separated
\end{lstlisting}

\begin{lstlisting}[basicstyle=\ttfamily\tiny,breaklines=true,
  breakatwhitespace=true,keepspaces=true,columns=fullflexible]
lambda_selection:
  lambda_delta: {value: 1.0, fixed: true, swept: false}
  lambda_pde:
    source_arm: standard_pinn
    candidates: [0.0, 0.01, 0.1, 1.0]
  lambda_gamma:
    source_arm: rung3_delta_gamma_vega
    candidates: [0.3, 1.0, 3.0]
  lambda_vega:
    source_arm: rung3_delta_gamma_vega
    candidates: [0.3, 1.0, 3.0]
  tune_on: validation_only
  robustness_row:
    description: "also report the confirmatory contrast at the rung3-sourced lambda_pde"
    status: robustness_result_not_a_second_confirmatory_test
\end{lstlisting}

The selected weights were
$(\lambda_{\mathrm{pde}},\lambda_D,\lambda_G,\lambda_v)
=(0.01,1,1,1)$. The registered alternative sourced from rung 3 selected
$\lambda_{\mathrm{pde}}=0$ and was reported as a robustness row.

\paragraph{Amendment record.} The contract evolved through four documented,
pre-run amendments. All preceded training and model evaluation; amendment 3 used a
production-scale \emph{protocol diagnostic} to correct the clipping rationale, but
no trained-arm outcome.

\begin{table}[h]
\caption{Contract amendments before the frozen label artifact. ``Pre-eval'' means
before any trained-arm evaluation.}
\label{tab:amendments-s2s}
\centering\footnotesize\setlength{\tabcolsep}{4pt}
\begin{tabular}{@{}lll p{0.52\linewidth}@{}}
\toprule
ID & Date & Pre-eval & Reason and effect \\
\midrule
A1 & 2026-07-31 & yes & Declared formerly implicit thresholds; recorded the realised 43-step clock; distinguished ``monotonicity not demonstrated'' from flatness; fixed validation-only weight sourcing. \\
A2 & 2026-08-04 & yes & Made the $10\%$ confirmatory threshold numeric; separated absent from broken evaluations in the verdict vocabulary; declared the corruption-noise ladder of the pre-training detectability diagnostic, and its region of validity. \\
A3 & 2026-08-05 & yes & Corrected the measured clipping direction in that same diagnostic, demoted its $0.20$ noise rung to diagnostic-only, and fixed the extension contingency. The triggering data were clip measurements, not arm outcomes. \\
A4 & 2026-08-05 & yes & Declared the $20\%$ parameter-point validation split, $64$ state--contract--time draws per parameter point, and full Monte Carlo coverage. \\
\bottomrule
\end{tabular}
\end{table}

\begin{table}[!htbp]
\caption{Registered endpoints and their outcomes. ``Not demonstrated'' is the
contract's own wording for a dose-response that fails its monotonicity criterion,
which is not the claim that the response is flat. The in-class check is declared
\texttt{NOT\_PASS\_FAIL}: it flags a registered expectation motivated by the similar
pricing and Greek accuracy of the supervised models in \citet{sakuma2026a},
not a hedging null established by that paper. Intervals below are in loss-difference units;
the dose probability is a one-sided seed-bootstrap tail probability.}
\label{tab:verdicts}
\centering\footnotesize\setlength{\tabcolsep}{4pt}
\begin{tabular}{@{}p{0.26\linewidth}p{0.20\linewidth}p{0.28\linewidth}p{0.14\linewidth}@{}}
\toprule
Endpoint & Declared & Observed & Verdict \\
\midrule
Confirmatory: misspecified $J$, rung 3 vs baseline & $\ge 10\%$ relative, CI excludes 0 & $+0.02\%$, CI $[-0.0242,+0.0242]$ & fail \\
Order attribution: rung 2 beats rung 1 & CI excludes 0 & $+0.0574$, CI $[+0.0527,+0.0621]$; $-1.2\%$ relative & fail (reversal) \\
Dose-response monotonicity & $\rho_S>0$ and $P(\rho_S\le 0)<0.05$ & $\rho_S = 0.029$, $P = 0.074$ & not demonstrated \\
Held-out $\partial_{SS}u$ and $\partial_v u$ RMSE at value parity & $\ge 15\%$ each, value parity within $10\%$ & reductions $0.864$--$0.930$; value improved $82$--$91\%$ & fail (on parity) \\
In-class consistency check & not pass/fail & $+27.2\%$ gap at zero cost & flag \\
Mechanism rule & zero-cost gap or cost/turnover pattern & zero-cost gap; later controls do not isolate a Gamma-specific cause & channel (i), qualified \\
Departure sweep: locate a favourable regime & --- & none at the $1\%$ tier & no decisive regime \\
\bottomrule
\end{tabular}
\end{table}

% =============================================================================
\section{Reference construction and per-derivative solver spread}
\label{app:oracle}
Section~\ref{sec:setup} reports agreement across available numerical legs.
For derivatives, the measured spread is a diagnostic scale for arm comparisons;
it is not a bound on reference error or a guarantee of an ordering.

\paragraph{Construction, and which legs run where.} Leg~1 analytically differentiates the
trap-free Heston $P_1/P_2$ Fourier integrands
\citep{heston1993closed,albrecher2007little} before Gauss--Legendre quadrature. Leg~2 differentiates a Fang--Oosterlee cosine-expansion
pricer \citep{fang2008novel} with fourth-order central stencils and a three-step
Richardson sweep. It separately implements the same $\phi_2$ formula as leg~1, so their agreement
cannot exclude a shared formula error. Leg~3 is
Monte Carlo, pathwise for the first derivative and mixed pathwise/likelihood-ratio
for curvature, with antithetic pairs and a control variate; it returns standard
errors and has no estimator for the mixed derivative. Leg~4 is a Craig--Sneyd ADI
solve \citep{craig1988alternating}, routed only where the variance process is
stressed. The interpolated anchor (Feller ratio $1.04$) draws all four legs and a
four-fold Monte-Carlo path multiplier, $800{,}000$ paths; the extrapolated anchor
(Feller ratio $1.78$) draws three, $200{,}000$ paths, and has no ADI leg. The legs have distinct numerical procedures but do not provide four fully
independent certifications.

Per grid point the reference value is the median across legs and the reported
uncertainty is the largest pairwise absolute disagreement, folded with the
Monte-Carlo standard error; a point whose largest disagreement exceeds
$10^{-3}$ times a wing-floored scale, plus $3\sigma_{\mathrm{MC}}$ for any pair
involving the stochastic leg, is masked. The rule was declared before training.

\paragraph{Deterministic spread and stochastic uncertainty.}
Table~\ref{tab:spread} reports RMS disagreement among deterministic legs alongside
the broader envelope. At the extrapolated anchor the two characteristic-function
legs agree to $1.6\times10^{-8}$ on $\partial_Su$ and $3.3\times10^{-8}$ on
$\partial_{SS}u$ in RMS. This agreement cannot rule out their shared bias.
The curvature envelope is $2.65$ times larger there despite the interpolated
anchor's lower Feller ratio. The stored Monte-Carlo standard errors have ratio
$2.63$; a nominal $n^{-1/2}$ adjustment for the fourfold path-budget difference
reduces this to $1.32$. This is a scaling diagnostic, not an equal-budget rerun or
a bias/convergence certificate. The earlier pointwise MC-dominance percentages
are omitted because an audit did not consistently reproduce them under explicit
dominance definitions.

\begin{table}[!htbp]
\caption{Per-derivative reference uncertainty at the two evaluation anchors, in the
units of Table~\ref{tab:rmse} (curvature $\times 10^{-3}$), as root mean square over
the retained points each arm is scored on. ``Envelope'' is the quantity the paper
reports: largest pairwise leg disagreement folded with the Monte-Carlo standard
error. ``Deterministic'' is the largest pairwise disagreement among the
non-stochastic legs alone, recomputed on the frozen grids. The ratio is envelope
divided by deterministic spread; neither quantity bounds common solver bias.}
\label{tab:spread}
\centering\footnotesize\setlength{\tabcolsep}{5pt}
\begin{tabular}{@{}llrrr@{}}
\toprule
Regime & Derivative & envelope & deterministic & ratio \\
\midrule
interpolated & $\partial u/\partial S$ & $0.000276$ & $6.96\times10^{-6}$ & $40$ \\
(4 legs, $8\times10^{5}$ paths) & $\partial^{2}u/\partial S^{2}$ & $0.768$ & $0.00156$ & $491$ \\
 & $\partial u/\partial v$ & $0.1083$ & $0.000656$ & $165$ \\
\midrule
extrapolated & $\partial u/\partial S$ & $0.000397$ & $1.61\times10^{-8}$ & $2.5\times10^{4}$ \\
(3 legs, $2\times10^{5}$ paths) & $\partial^{2}u/\partial S^{2}$ & $2.035$ & $3.31\times10^{-5}$ & $6.1\times10^{4}$ \\
 & $\partial u/\partial v$ & $0.1270$ & $5.07\times10^{-6}$ & $2.5\times10^{4}$ \\
\bottomrule
\end{tabular}
\end{table}

\paragraph{Limits of the resolution heuristic.} The reference is a median across
available legs. Its agreement with characteristic-function values does not prove
truth or independence. Even a valid RMS reference-error bound $\epsilon$ would
allow a difference of two RMSEs to move by up to $2\epsilon$; measured solver
spread need not be such a bound. Exceeding one envelope therefore supplies a
sensitivity diagnostic, not a one-sided guarantee. Falling below it is likewise
not evidence of equality.

\paragraph{Which contrasts in Table~\ref{tab:rmse} survive.} Table~\ref{tab:resolve}
divides each arm-to-arm RMSE difference by the envelope for that derivative and
anchor. The convention follows the study's registered
\texttt{improvement\_to\_oracle\_noise} statistic; it is a resolution heuristic, not
a hypothesis test, and it ignores seed dispersion, which is reported separately in
Table~\ref{tab:rmse}'s source.

\begin{table}[!htbp]
\caption{Arm-difference divided by the reference envelope of Table~\ref{tab:spread}.
Positive means the second arm is more accurate. A magnitude below $1$ falls below the declared
resolution heuristic; this is not an inferential significance threshold.}
\label{tab:resolve}
\centering\footnotesize\setlength{\tabcolsep}{5pt}
\begin{tabular}{@{}lrrrrrr@{}}
\toprule
& \multicolumn{3}{c}{interpolated} & \multicolumn{3}{c}{extrapolated} \\
\cmidrule(lr){2-4}\cmidrule(lr){5-7}
Contrast & $\partial_S u$ & $\partial_{SS} u$ & $\partial_v u$
         & $\partial_S u$ & $\partial_{SS} u$ & $\partial_v u$ \\
\midrule
baseline $\to$ rung 1 & $156.2$ & $6.9$ & $67.1$ & $86.8$ & $2.2$ & $45.0$ \\
baseline $\to$ rung 3 & $184.7$ & $9.1$ & $85.0$ & $97.3$ & $2.6$ & $54.5$ \\
rung 1 $\to$ rung 2 & $26.2$ & $2.0$ & $9.9$ & $11.4$ & $\mathit{0.5}$ & $3.5$ \\
rung 2 $\to$ rung 3 & $2.2$ & $\mathit{0.2}$ & $8.0$ & $\mathit{-0.97}$ & $\mathit{-0.1}$ & $6.0$ \\
\bottomrule
\end{tabular}
\end{table}

Three readings follow, and Sections~\ref{sec:survives}--\ref{sec:budget} are written
to them.

\emph{The headline survives everywhere.} Every baseline-to-supervised contrast
exceeds the envelope at both anchors, by factors from $2.2$ to $185$. The $82$ to
$93\%$ reductions exceed this diagnostic scale, conditional on the reference.

\emph{The curvature rung is resolved at one anchor only.} The rung-1-to-rung-2
curvature gap is $2.0$ envelopes at the interpolated hole and $0.5$ at the
extrapolated one. Section~\ref{sec:curvature} claims the increment at the
interpolated anchor and declines to read its extrapolated magnitude finely. The
transfer result of Section~\ref{sec:transfer} lives in the $\partial_v u$ column and
is resolved at both anchors, at $9.9$ and $3.5$ envelopes.

\emph{The two extrapolated reversals are not resolved, and the appendix says so
plainly.} The flagged risk was that Section~\ref{sec:budget}'s non-monotonicity
rests on $0.641 \to 0.823$ on $\partial^{2}u/\partial S^{2}$ and
$0.003399 \to 0.003784$ on $\partial u/\partial S$. Those gaps are $0.09$ and $0.97$
envelopes. Both sit inside the reference uncertainty, so this benchmark cannot call
either a degradation, and Section~\ref{sec:budget} does not. What survives there is
the resolved gain on the derivative rung 3 adds, at $8.0$ and $6.0$ envelopes.

A sharper form of the same limit: at the extrapolated anchor the curvature RMSE of
rungs 1, 2 and 3 is $0.77$, $0.31$ and $0.40$ times the envelope. All three errors fall below the envelope, so fine orderings remain conditional
on the reference rather than independently certified. Only the gap from the unsupervised arms, which are $3.0$ and $3.5$ envelopes
away, is.

\paragraph{Masking: rate, location, and one correction.} The union mask over all
five reference quantities removes $17.11\%$ of the interpolated grid
($3{,}704$ of $21{,}648$ points) and $9.23\%$ of the extrapolated grid
($1{,}999$ of $21{,}648$). Per quantity the interpolated rates are $0.008$ (value),
$0.038$, $0.073$, $0.080$ and $0.116$ for the first derivative, curvature, second
state and mixed derivative; the extrapolated rates are $0.016$, $0.014$, $0.055$,
$0.044$ and $0.052$. Table~\ref{tab:rmse} scores every arm on the same union mask.

The mask is a relative-tolerance artefact and sits where that tolerance is hardest
to meet. It concentrates overwhelmingly at the shortest maturity, $94.3\%$ of the
$\tau = 0.04$ slice at the interpolated anchor and $78.7\%$ at the extrapolated one,
against $0.077$ to $0.102$ beyond $\tau = 0.45$ at the interpolated anchor and
$0.000$ to $0.020$ beyond $\tau = 0.29$ at the extrapolated one; and it peaks just out
of the money rather than deep out of it, $47.5\%$ in the $S/K \in [0.75, 0.9)$ band
at the interpolated anchor against $18.1\%$ below $0.75$ and $6.8\%$ above $1.3$. In
curvature the pattern is U-shaped, heaviest in the lowest decile, where a relative
tolerance on $|\partial^{2}u/\partial S^{2}| \le 3\times10^{-12}$ demands absolute
agreement no solver delivers, and in the highest.

The extrapolated anchor is less masked, $9.23\%$ versus $17.11\%$. The fourth
leg at the interpolated anchor introduces additional pairwise checks. These masks
are selective; shared support makes arm comparisons consistent on retained points
but does not establish performance over excluded states.

% =============================================================================
\section{Training configuration and excision protocol}
\label{app:config}
The main ladder and baseline arms share architecture, optimiser, data and step
budget; their loss terms differ as listed below. The extra-value-data sweep is an
exception. Table~\ref{tab:trainconfig} states the shared
budget once and Table~\ref{tab:arms-s2s} the per-arm difference, so that the
identical-configuration claim of Section~\ref{sec:setup} is checkable rather than
asserted.

\begin{table}[!htbp]
\caption{Shared configuration. Every entry is identical across the ladder and
baseline arms; the extra-value-data control of Table~\ref{tab:arms-s2s} departs from
it in the weight, row-budget and selection settings described below.}
\label{tab:trainconfig}
\centering\footnotesize\setlength{\tabcolsep}{5pt}
\begin{tabular}{@{}lp{0.64\linewidth}@{}}
\toprule
Component & Setting \\
\midrule
Inputs & $(S, K, \tau, \kappa, \theta, \xi, \rho, v_0)$, affinely mapped to $[-1,1]$ by fixed contract ranges inside the forward pass with no clamp, so evaluated regimes outside those ranges extrapolate rather than saturate \\
Network & $4 \times 64$ \texttt{tanh} multilayer perceptron, $13{,}121$ parameters, \texttt{float32}; no hard-constrained ansatz, terminal treatment identical across arms \\
Optimiser & Adam, learning rate $3\times10^{-3}$ cosine-annealed to $3\times10^{-5}$, gradient-norm clip $1.0$ \\
Budget & $20{,}000$ steps, minibatch $1{,}024$ label rows, $4{,}096$ residual collocation points redrawn each epoch ($834$ epochs); fixed-step mode (called matched-epoch in code), early stopping disabled \\
Labels & $23{,}623$ training and $5{,}936$ validation rows from the frozen artifact \texttt{v6-labels-20260812}: $512$ coefficient draws crossed with $64$ shared $(S,K,\tau)$ triples, mask rate $0.098$ \\
Loss normalisation & supervised terms divided by second moments computed once over the retained training split; the Gamma scale uses true consensus \texttt{gamma\_ref}, shared across arms. Without it the curvature term sits about six orders of magnitude below the value term at equal relative error and is numerically inert \\
Weights & $(\lambda_{\mathrm{pde}}, \lambda_1, \lambda_2, \lambda_v) = (10^{-2}, 1, 1, 1)$, selected on validation only, sourced per Appendix~\ref{app:thresholds} \\
Seeds & $5$ per arm; $10$ on the ladder arms and the residual-only baseline \\
Hardware & CPU, single process; per-arm wall clock and peak memory in Table~\ref{tab:arms-s2s}. Processor model and a project-wide compute total were not recorded \\
\bottomrule
\end{tabular}
\end{table}

\begin{table}[!htbp]
\caption{Arms: label flags and measured cost. \emph{res.} is the physics residual;
$\partial_S$, $\partial_{SS}$, $\partial_v$ are the three label terms. Wall clock is
the per-seed mean; peak memory is the seed-42 measurement. Derivative evaluations
count first- and second-order AD calls, rather than only double-backward passes.
$\ast$ The extra-value-data control uses residual weight $1$ instead of $10^{-2}$
and $12{,}288$ retained rows instead of $23{,}623$. Its checkpoint and run log
have conflicting configurations; its cost figures are withheld.}
\label{tab:arms-s2s}
\centering\footnotesize\setlength{\tabcolsep}{4pt}
\begin{tabular}{@{}lccccrrrr@{}}
\toprule
Arm & res. & $\partial_S$ & $\partial_{SS}$ & $\partial_v$ & seeds & wall (s) & deriv.\ evals & peak (MB) \\
\midrule
value-only feedforward & -- & -- & -- & -- & $5$ & $17$ & $0$ & $383$ \\
residual-only baseline & \checkmark & -- & -- & -- & $10$ & $377$ & $60{,}120$ & $978$ \\
rung 1, $+\partial_S u$ & \checkmark & \checkmark & -- & -- & $10$ & $414$ & $80{,}160$ & $941$ \\
rung 2, $+\partial_{SS} u$ & \checkmark & \checkmark & \checkmark & -- & $10$ & $471$ & $100{,}200$ & $1009$ \\
rung 3, $+\partial_v u$ & \checkmark & \checkmark & \checkmark & \checkmark & $10$ & $451$ & $100{,}200$ & $1043$ \\
supervised, residual-free & -- & \checkmark & \checkmark & \checkmark & $5$ & $64$ & $40{,}080$ & $695$ \\
\midrule
curvature $L^{2}$ penalty & \checkmark & -- & -- & -- & $5$ & $465$ & $100{,}200$ & $1267$ \\
shuffled labels & \checkmark & \checkmark & \checkmark & \checkmark & $5$ & $433$ & $100{,}200$ & $1095$ \\
extra value data$^{\ast}$ & \checkmark & -- & -- & -- & $5$ & --- & --- & --- \\
\bottomrule
\end{tabular}
\end{table}

\paragraph{Excision.} Training coefficients are a seeded Latin hypercube on
$\kappa \in [1,4]$, $\theta \in [0.02,0.12]$, $\xi \in [0.20,0.60]$,
$\rho \in [-0.80,-0.20]$, $v_0 \in [0.01,0.12]$, rejecting any draw with
$2\kappa\theta/\xi^{2} < 0.40$ (an oracle-trust cutoff, not the Feller condition).
Around each of the five named regimes a ball is excised: in coordinates
range-normalised to the unit cube, every draw within Euclidean distance
$0.10\sqrt{5} = 0.2236$ of an anchor, that is $10\%$ of the cube diagonal, is
rejected. Realised: no sampled point falls inside any ball, the nearest lying at
$1.03$ to $1.64$ times the radius, minimum normalised anchor distance $0.2310$ to
$0.3670$; minimum realised Feller ratio $0.4061$.

The validation split is $20\%$ of \emph{coefficient points}, not of rows, so a
held-out regime cannot leak into validation through a shared coefficient draw. Both
evaluated regimes are therefore absent from training and from every selection score.
The interpolated anchor
$(\kappa, \theta, \xi, \rho, v_0) = (1.5, 0.04, 0.34, -0.70, 0.02)$ sits inside the
sampled box in every coordinate and is a hole punched in the support. The
extrapolated anchor $(2.0, 0.04, 0.30, -0.90, 0.04)$ has $\rho = -0.90$ against a
sampled range whose lower bound is $-0.80$ and whose minimum realised draw is
$-0.7989$: the ball intersects the cube only where it overlaps, and the anchor
itself lies outside the sampled support entirely, so that regime is extrapolation
and not interpolation into a hole.

The contract grid is $41 \times 33 \times 16$ over $S \in [50,150]$,
$K \in [60,140]$, $\tau \in [0.04,1]$ at $r = 0.02$, $q = 0$. Training labels use
$64$ shared $(S,K,\tau)$ triples, giving sparse local coverage near the hedged
contract. Absence of a nearby sampled triple does not by itself establish
extrapolation outside the training support in these coordinates.

\paragraph{Departure variants.} The Bates generating process adds a
compound-Poisson component \citep{bates1996jumps} to the baseline anchor
$(\kappa, \theta, \xi, \rho, v_0) = (2.0, 0.04, 0.30, -0.50, 0.04)$, over
$\lambda_j \in \{0, 0.1, 0.25, 0.5\}$ and $\sigma_j \in \{0.05, 0.10, 0.15\}$ at
$\mu_j = -0.10$. The Merton variant \citep{merton1976option} replaces the
two-state diffusion with a one-state jump-diffusion at $\sigma = 0.20$,
$\lambda_j = 0.25$, $\mu_j = -0.10$, $\sigma_j = 0.10$, so the surrogate stays
conditioned on a second state the world does not have. Every arm is handed the
baseline anchor's coefficients at decision time, never the true ones. The
$\lambda_j = 0$ cell is the contract's in-class consistency check and reproduces the
Heston cell exactly, $1.7371$ against $1.7371$. It is the least-departed cell in
Table~\ref{tab:severity} and carries the largest advantage in it, $+27.7\%$ at zero
cost. A second registered check compares the same corner with an expectation of
near-indifference motivated by the similar pricing and Greek accuracy of the three
supervised models in \citet{sakuma2026a}. It \emph{flags}: the gap exceeds that
registered expectation. Their accuracy comparison does not itself predict a
hedging null in our setting. We retain the flag as a declared consistency check,
not a pass/fail endpoint.

\paragraph{The registered severity sweep, read out.} Table~\ref{tab:severity} gives
the sweep Section~\ref{sec:survives} summarises. At zero action cost the advantage
of the fully supervised arm over the residual-only baseline falls monotonically in
\emph{each} departure coordinate, all six one-dimensional chains decreasing, from
$+27.7\%$ at the in-class cell to $-4.2\%$ at the most departed one, changing sign
between $\lambda_j = 0.5,\ \sigma_j = 0.05$ and $\lambda_j = 0.5,\ \sigma_j = 0.10$.
Every pooled interval excludes zero except the sign-crossing cell itself. At the
registered $1\%$ cost the whole table is harm and no cell is favourable, which is
what produced the registered \texttt{no\_decisive\_regime} verdict: that verdict
reads only the $1\%$ tier. The ordering is ordinal. We know which cells are more
departed, and the advantage decays in that order; we did not compute a distance from
the assumed model class, so the axis has no calibrated scale and the table is not a
severity-versus-advantage \emph{curve} in the metric sense.

\begin{table}[!htbp]
\caption{Registered departure sweep, delta-only $J$ (lower is better), five seeds,
shared trajectories. ``Advantage'' is the relative improvement of rung 3 over the
residual-only baseline; positive favours supervision. $\dagger$ marks the one cell
whose pooled paired-bootstrap interval on the difference includes zero. The
$\lambda_j = 0$ row is the in-class consistency check, not a departure.}
\label{tab:severity}
\centering\footnotesize\setlength{\tabcolsep}{5pt}
\begin{tabular}{@{}llrrrr@{}}
\toprule
& & \multicolumn{2}{c}{zero cost} & \multicolumn{2}{c}{$1\%$ cost} \\
\cmidrule(lr){3-4}\cmidrule(lr){5-6}
$\lambda_j$ & $\sigma_j$ & rung 3 & advantage & rung 3 & advantage \\
\midrule
$0$ (in class) & $0.05$ & $1.737$ & $+27.7\%$ & $4.818$ & $-6.9\%$ \\
\midrule
$0.10$ & $0.05$ & $2.096$ & $+20.5\%$ & $5.082$ & $-7.8\%$ \\
$0.10$ & $0.10$ & $2.358$ & $+17.9\%$ & $5.330$ & $-7.5\%$ \\
$0.10$ & $0.15$ & $2.699$ & $+15.6\%$ & $5.669$ & $-7.0\%$ \\
$0.25$ & $0.05$ & $2.723$ & $+10.6\%$ & $5.533$ & $-9.1\%$ \\
$0.25$ & $0.10$ & $3.455$ & $+7.0\%$ & $6.207$ & $-8.1\%$ \\
$0.25$ & $0.15$ & $4.384$ & $+4.9\%$ & $7.113$ & $-6.5\%$ \\
$0.50$ & $0.05$ & $3.582$ & $+0.8\%^{\dagger}$ & $6.205$ & $-11.1\%$ \\
$0.50$ & $0.10$ & $5.057$ & $-3.0\%$ & $7.532$ & $-9.6\%$ \\
$0.50$ & $0.15$ & $6.870$ & $-4.2\%$ & $9.280$ & $-7.8\%$ \\
\midrule
state removed & & $2.925$ & $+9.8\%$ & $5.845$ & $-3.7\%$ \\
\bottomrule
\end{tabular}
\end{table}

% =============================================================================
\section{Supporting exhibits}
\label{app:exhibits}

\paragraph{Value accuracy does not imply derivative accuracy.} Table~\ref{tab:motiv}
is the exhibit behind Section~\ref{sec:intro}'s opening clause. Both unsupervised
arms match the reference solution to a few percent while their curvature is wrong by
tens of percent, and the ordering on the value is not the ordering on the curvature:
the residual-only arm has the better value fit and the worse curvature.

\begin{table}[!htbp]
\caption{Relative RMSE against the reference, at the two held-out coefficient
regimes, five seeds. Controlling the value in $L^{2}$ leaves curvature error one to
two orders of magnitude larger in relative terms.}
\label{tab:motiv}
\centering\footnotesize\setlength{\tabcolsep}{6pt}
\begin{tabular}{@{}lrrrr@{}}
\toprule
& \multicolumn{2}{c}{interpolated} & \multicolumn{2}{c}{extrapolated} \\
\cmidrule(lr){2-3}\cmidrule(lr){4-5}
Arm & $u$ & $\partial^{2}u/\partial S^{2}$ & $u$ & $\partial^{2}u/\partial S^{2}$ \\
\midrule
residual-only baseline & $0.0155$ & $0.4096$ & $0.0128$ & $0.3505$ \\
value-only feedforward & $0.0328$ & $0.2664$ & $0.0318$ & $0.2206$ \\
rung 3 (fully supervised) & $0.0014$ & $0.0446$ & $0.0022$ & $0.0498$ \\
\bottomrule
\end{tabular}
\end{table}

\paragraph{The information-free controls.} Table~\ref{tab:controls} gives the three
controls of Section~\ref{sec:controls} with seeds and dispersion, together with the
first-derivative column, which is what separates ``no signal'' from ``objective
conflict'': the curvature penalty is worse than the baseline on the derivative it
never touches as well as on the one it shapes. The two configuration-matched
controls are the ones the argument rests on; the third is discussed below.

\paragraph{Scope of the extra-value-data sweep.} The saturation driver uses
$\lambda_{\mathrm{pde}}=1$ rather than the selected $10^{-2}$. Its nominal base
$N=4{,}096$ is also smaller than the main arms' actual $23{,}623$ retained rows:
even $5N=20{,}480$ remains below the main budget. The selected $3N$ checkpoint uses
$12{,}288$ rows. Thus the level comparison in Table~\ref{tab:controls} is not
matched in weight or data.

Within the sweep, validation curvature relative RMSE changes from
$0.647\pm0.095$ to $0.596\pm0.015$ across $1N$--$5N$. The declared $2\%$ rule
selects $3N$ ($0.605\pm0.017$); its doubled-width control gives
$0.709\pm0.061$. This describes the tested protocol and does not rule out
optimization limitations. The sweep uses its own validation-selection protocol;
its scores are distinct from held-out RMSE. A substitution claim would require
matched actual row budgets, residual weights and selection rules. The penalty and
shuffled-label controls are configuration-matched, but their failures do not prove
a general information bound.

\begin{table}[!htbp]
\caption{Control arms, held-out relative RMSE, mean $\pm$ seed standard deviation
over five seeds. $\lambda_{\mathrm{pde}}$ and $N$ are the residual weight and
actual retained price-row count. Curvature relative RMSE uses the declared central
moneyness support; absolute RMSE in Table~\ref{tab:rmse} uses the full retained grid.
The extra-data row is not a matched comparison.}
\label{tab:controls}
\centering\footnotesize\setlength{\tabcolsep}{4pt}
\begin{tabular}{@{}lrrrrr@{}}
\toprule
& & & \multicolumn{2}{c}{$\partial^{2}u/\partial S^{2}$ rel.\ RMSE} & $\partial u/\partial S$ \\
\cmidrule(lr){4-5}
Arm & $\lambda_{\mathrm{pde}}$ & $N$ & interpolated & extrapolated & interpolated \\
\midrule
residual-only baseline & $0.01$ & $23{,}623$ & $0.4096 \pm 0.0106$ & $0.3505 \pm 0.0158$ & $0.0749$ \\
curvature $L^{2}$ penalty, no labels & $0.01$ & $23{,}623$ & $0.9822 \pm 0.0068$ & $0.9776 \pm 0.0079$ & $0.5389$ \\
shuffled labels & $0.01$ & $23{,}623$ & $0.4075 \pm 0.0343$ & $0.3666 \pm 0.0300$ & $0.1360$ \\
extra value data (plateau) & $1.00$ & $12{,}288$ & $0.6986 \pm 0.0118$ & $0.6147 \pm 0.0116$ & $0.2108$ \\
rung 3 supervision & $0.01$ & $23{,}623$ & $\mathbf{0.0446 \pm 0.0028}$ & $\mathbf{0.0498 \pm 0.0018}$ & $\mathbf{0.0038}$ \\
\bottomrule
\end{tabular}
\end{table}

\paragraph{Reconstructing the variance-minimising controller.} The combination
$\partial_S u + (\rho\xi/S)\,\partial_v u$ is the local-risk-minimising ratio for
this state \citep{follmer1991hedging}. Any arm carrying both first derivatives can
assemble it from its own outputs with no further reference, so Table~\ref{tab:mv-s2s}
is the point at which the field result and the decision result meet. It is
exploratory and post-registered, moves no registered verdict, and reuses the frozen
checkpoints and the same shared trajectories every other cell faces; $\rho$ and $\xi$
in the correction are the assumed-model values, so no true-process parameter enters.
Point estimates only: these are five-seed means against a pre-registered $2\%$ floor,
and we claim no interval for them.

Both arms carrying curvature labels land inside the floor, as does the residual-free
supervised arm; the two arms without curvature labels miss it by $3.3$ and $3.4$
percent, and the residual-damaged baseline by $42.8$. The reference the registered
metric scores against, the exact first derivative alone, is itself $24.7\%$ worse
than this controller, which is the sense in which the registered endpoint could
not have rewarded accuracy at the accuracy supervision reaches.

\begin{table}[!htbp]
\caption{Delta-only $J$ (lower is better) when the controller is built from each
arm's own two first derivatives, on the confirmatory cell's combined perturbation
at zero action cost, five seeds,
shared trajectories. The gap is relative to the analytic construction; negative
beats it. The pre-registered negligibility floor is $2\%$.}
\label{tab:mv-s2s}
\centering\footnotesize\setlength{\tabcolsep}{6pt}
\begin{tabular}{@{}lrrl@{}}
\toprule
Fields driving the controller & $J$ & gap & within $2\%$ floor \\
\midrule
analytic construction & $1.6696$ & --- & reference \\
rung 2 ($+\partial_{SS}u$) & $\mathbf{1.6563}$ & $-0.79\%$ & yes \\
rung 3 ($+\partial_v u$) & $1.6700$ & $+0.03\%$ & yes \\
supervised, residual-free & $1.6793$ & $+0.58\%$ & yes \\
rung 1 ($+\partial_S u$) & $1.7240$ & $+3.26\%$ & no \\
value-only feedforward & $1.7261$ & $+3.38\%$ & no \\
residual-only baseline & $2.3843$ & $+42.81\%$ & no \\
\midrule
exact first derivative alone & $2.0816$ & $+24.68\%$ & no \\
\bottomrule
\end{tabular}
\end{table}

\paragraph{Delta accuracy on the full grid and the hedging region.}
The registered accuracy metric uses the full grid, $S\in[50,150]$,
$K\in[60,140]$, $\tau\in[0.04,1]$. The hedge trades one contract, with $K=100$
and remaining expiry $\tau\in[0.08,0.25]$. Table~\ref{tab:region} restricts the
same metric to the corresponding strike, expiry interval and the contract's
predefined central range $S/K\in[0.75,1.30]$. It is a restricted-grid comparison; it does not
weight states by how frequently trading trajectories visit them.

As summarised in Section~\ref{sec:coda}, delta accuracy reverses the two models
without derivative labels: the residual-only model is better on the full grid and
worse in the restricted region, in all five seeds at both regimes. Gamma also
reverses in the mean, but its full-grid ordering holds in only three of five and
four of five seeds. We therefore highlight the delta result in the main text.
Across training labels in the same expiry and $S/K$ ranges, pooling strikes, the
mask rate is $1.30$ times the rate outside. This indicates thinner reference-label
coverage in that part of the domain.

\begin{table}[!htbp]
\caption{Delta RMSE for the two models without derivative labels, on the full
parameter grid and the restricted hedging region, five seeds. Lower is better.
The rankings reverse; the regional metric is not trajectory-weighted.}
\label{tab:region}
\centering\footnotesize\setlength{\tabcolsep}{6pt}
\begin{tabular}{@{}lrrrr@{}}
\toprule
& \multicolumn{2}{c}{interpolated} & \multicolumn{2}{c}{extrapolated} \\
\cmidrule(lr){2-3}\cmidrule(lr){4-5}
Arm & full grid & region & full grid & region \\
\midrule
residual-only baseline & $\mathbf{0.0537}$ & $0.0937$ & $\mathbf{0.0424}$ & $0.0689$ \\
value-only feedforward & $0.0881$ & $\mathbf{0.0409}$ & $0.0815$ & $\mathbf{0.0227}$ \\
\bottomrule
\end{tabular}
\end{table}

\end{document}
